\chapter{基础设施层：算力、芯片与 AI 工厂}
\label{ch:infrastructure}

% The foundational layer enabling all AI applications.
% Covers GPU cloud, custom silicon, training clusters, inference serving,
% networking, and edge compute.

% §2.1 GPU Cloud and Compute-as-a-Service
% =============================================================
% §2.1  GPU Cloud and Compute-as-a-Service
% Book: The AI Startup Landscape
% =============================================================

\section{GPU 云与算力即服务}
\label{sec:infra-gpu-cloud}

想象这样一个场景：2022年末，一座位于新泽西州的加密货币矿场正在为
以太坊转向权益证明（Proof-of-Stake）而头疼——数千块高端GPU突然失去
了工作。矿场老板面临两个选择：以废铁价甩卖硬件，或者赌一把全新的市场。
他们选择了后者。不到三年，这座矿场变成了市值超过400亿美元的上市公司，
客户名单上写满了微软、Meta和OpenAI的名字。

这就是CoreWeave的故事，也是整个GPU云（Neocloud）赛道的缩影。当
ChatGPT点燃了全球对大模型的狂热，算力突然成为比石油更稀缺的资源。
传统云巨头——AWS、Azure、GCP——虽然拥有海量基础设施，却在GPU供应
上面临数月的排队周期。一群嗅觉灵敏的创业者迅速填补了这个缺口，用
各种方式把GPU算力送到饥渴的AI团队手中：有人建造巨型数据中心，有人
做P2P市场，有人用天然气废气发电，有人把整个流程抽象成一行Python代码。

到2026年，GPUaaS（GPU即服务）市场规模已达73.6亿美元，预计2031年
将增长至264亿美元，复合年增长率29.1\%（Mordor Intelligence, 2026）。
全球已涌现100多家Neocloud创业公司，共同构成了AI基础设施层中最资本密集、
增长最快、竞争最激烈的赛道。

\begin{keybox}[GPU 云赛道：核心融资数据一览]
\small
\begin{tabularx}{\textwidth}{l X r r}
\toprule
\rowcolor{figLightGray}
\textbf{公司} & \textbf{最新融资} & \textbf{总融资/估值} & \textbf{2025营收} \\
\midrule
CoreWeave      & IPO 2025/3 (CRWV)          & 市值\$43.5B          & \$5.13B \\
Lambda Labs    & Series E \$1.5B (2025/11)  & 总融\$2.3B           & $\sim$\$505M \\
Crusoe Energy  & Series E \$1.375B (2025/10)& 估值\$10B+           & 未披露 \\
Nebius         & NASDAQ: NBIS               & 市值$\sim$\$23B      & \$529.8M \\
Together AI    & Series B \$305M (2025)     & \$3.3B$\to$\$7.5B   & ARR$\sim$\$1B \\
Modal          & Series B \$87M (2025/9)    & \$1.1B$\to$\$2.5B   & 未披露 \\
RunPod         & 仅融\$22M                   & ARR \$120M           & \$120M+ \\
Vast.ai        & 仅融\$4M                    & ---                  & 未披露 \\
SF Compute     & Series A \$40M (2025/12)   & 估值\$300M           & 未披露 \\
\bottomrule
\end{tabularx}

\medskip
\noindent 数据来源：Crunchbase, S-1 filing, 公司公告, TechCrunch,
The Information. 截至2026年3月。
\end{keybox}

% ===========================================================
\subsection{CoreWeave：从矿场到华尔街}
\label{subsec:coreweave}
% ===========================================================

CoreWeave的创始故事堪称AI时代最传奇的转型叙事。2017年，三位华尔街
交易员——Michael Intrator、Brian Venturo和Braden Perley——在新泽西
州创立了一家名为Atlantic Crypto的以太坊挖矿公司。彼时加密市场如火
如荼，GPU是印钞机的核心零件。

转折发生在2022年9月。以太坊完成"合并"（The Merge），从工作量证明
切换到权益证明，GPU矿机一夜之间变成了摆设。但这三位创始人看到了一
个更大的机遇：同样的NVIDIA GPU，在加密矿场中做哈希运算，换一套软件
栈就能做矩阵乘法——而后者正是大模型训练的核心操作。

\begin{whybox}[为什么加密矿场能转型AI？]
从技术角度看，加密挖矿和AI训练共享三个关键资源：\textbf{大规模GPU
集群}、\textbf{低成本电力}、以及\textbf{高密度散热能力}。矿场运营者
在过去五年中积累了以下AI数据中心同样需要的核心能力：
\begin{itemize}
  \item \textbf{GPU采购渠道}：与NVIDIA有成熟的大批量采购关系；
  \item \textbf{电力谈判}：擅长获取低价电力合约
    （通常\$0.03--0.05/kWh）；
  \item \textbf{物理基础设施}：已建成高功率密度的机房和冷却系统；
  \item \textbf{运维经验}：7\texttimes 24小时管理数千块GPU的经验
    直接可迁移。
\end{itemize}
不同之处在于，AI负载对\textbf{GPU间互联}（InfiniBand/NVLink）和
\textbf{网络拓扑}的要求远高于挖矿。CoreWeave的关键决策是大幅投资
网络基础设施，将松散的GPU矿机群升级为真正的HPC集群。
\end{whybox}

这一战略赌注收获了指数级回报。2023年，CoreWeave从Magnetar Capital
拿到了23亿美元的债务融资，随后又以NVIDIA提供的GPU作为抵押物获得更多
信贷额度。到2024年底，公司已签下微软作为锚定客户，合约金额据报道
超过百亿美元。

2025年3月28日，CoreWeave在纳斯达克上市（CRWV），IPO定价\$40，首日
微涨后一度跌破发行价。但随后股价一路攀升，在2025年6月触及\$187的
历史高点，市值一度逼近千亿美元。截至2026年3月，股价在\$80附近波动，
市值约435亿美元（CNBC, 2026年3月）——虽然从峰值回落了一半多，但这
依然是AI基础设施领域最大的纯GPU云上市公司。

CoreWeave 2025财年数据令人瞩目：营收51.3亿美元，同比增长730\%
（S-1 filing）。更惊人的是其\textbf{积压订单}（backlog）高达668亿
美元，意味着未来数年的收入已基本锁定。公司还以17亿美元收购了AI开发者
工具公司Weights \& Biases（W\&B），从纯基础设施向平台化演进——这一
动作表明，单纯卖算力的利润率终究会被压缩，CoreWeave需要沿价值链
向上攀升。

NVIDIA将CoreWeave列为"Exemplar Cloud Partner"，在GPU分配上给予
优先地位。这种深度绑定是一把双刃剑：它让CoreWeave在供给端获得了巨大
优势，但也意味着公司的命运高度依赖于与单一供应商的关系（参见本节末
的风险分析框）。对推理优化的技术细节感兴趣的读者，可参见
\textit{Emergence}第九章关于大规模推理系统的讨论。

% ===========================================================
\subsection{Lambda Labs：学术基因的超级算力云}
\label{subsec:lambda}
% ===========================================================

如果说CoreWeave是"矿工转型"，Lambda Labs则代表了另一条路径——从
深度学习工具起家，逐步扩展到全栈云服务。

Lambda由Stephen Balaban和Michael Balaban兄弟于2012年在硅谷创立，
最初是一家面向机器学习研究者的GPU工作站供应商。早年的客户以大学实验室
和研究机构为主，Lambda的品牌在学术圈建立了深厚的信任。随着大模型
时代来临，Lambda敏锐地将业务重心转向GPU云——不仅卖硬件，更提供按需
使用的云端A100/H100集群。

2025年11月，Lambda完成了15亿美元的E轮融资，由TWG Global和USIT领投，
累计融资总额达到23亿美元。公司将自己定位为"超级智能云"
（Superintelligence Cloud），目标是构建专为最大规模AI训练任务优化
的基础设施。2025年营收约5.05亿美元，虽然体量不及CoreWeave，但增速
同样惊人。

Lambda的差异化在于其\textbf{开发者体验}：公司保留了早期做工具的基因，
提供开箱即用的PyTorch/JAX环境、一键启动的Jupyter Hub、以及与主流
MLOps工具的深度集成。对于一个刚拿到NeurIPS论文审稿意见、急需跑大规模
实验的研究团队来说，Lambda的上手成本明显低于在AWS上从零配置。

% ===========================================================
\subsection{Crusoe Energy：用废气点燃AI}
\label{subsec:crusoe}
% ===========================================================

Crusoe Energy Systems的故事从一个环保主义者的洞察开始。2018年，
Chase Lochmiller和Cully Cavness在科罗拉多州注意到一个矛盾：油田
每天白白烧掉数百万立方英尺的伴生天然气（flare gas），而数据中心则在
疯狂寻找低成本电力。他们的解决方案简单而优雅——把移动数据中心搬到
油田旁边，直接用原本要烧掉的废气发电驱动GPU。

这一模式在ESG投资者中引起了强烈共鸣：减少甲烷排放的同时提供AI算力，
一举两得。2025年10月，Crusoe完成了13.75亿美元的E轮融资，估值超过
100亿美元。更引人注目的是，Crusoe是OpenAI/SoftBank发起的"Stargate"
超级计算项目的联盟成员，在德克萨斯州阿比林（Abilene）规划了一座
1.2GW的超大规模数据中心园区——一度被视为美国最大的AI基础设施项目
之一。

然而，2026年3月传来了变数：据报道，Oracle和OpenAI缩减了Abilene旗舰
Stargate站点的扩张计划，Meta正在NVIDIA的撮合下洽谈接手Crusoe的
部分产能（DCD, 2026年3月）。这一转折凸显了Neocloud赛道的一个核心
张力：\textbf{客户集中度风险}。当你的商业模式依赖于少数几个超大客户
的长期合约时，一个客户的战略调整就可能颠覆整个业务规划。

尽管如此，Crusoe在2025年DCD全球大奖中获得了"北美年度数据中心项目"
奖项，其废气发电模式作为差异化定位仍具独特价值。

% ===========================================================
\subsection{Nebius：从Yandex废墟中重生的AI云}
\label{subsec:nebius}
% ===========================================================

Nebius的故事是地缘政治重塑科技产业的活教材。2024年，俄罗斯搜索巨头
Yandex在制裁压力下被迫拆分，其海外AI基础设施业务以"Nebius"的名义
在阿姆斯特丹独立运营，由Yandex联合创始人Arkady Volozh领导。公司
继承了Yandex在大规模分布式系统方面的深厚技术积累，以及在芬兰和
欧洲各地已建成的数据中心。

Nebius以NASDAQ: NBIS在纳斯达克上市交易，2026年3月市值约230亿美元。
2025年营收5.298亿美元，同比暴增479\%。但真正让市场为之兴奋的是
2026年3月16日宣布的重磅消息：Nebius与Meta签署了一份五年期AI基础设施
供应协议，合同总额高达270亿美元（Reuters, 2026年3月16日）。消息
公布当天，NBIS股价跳涨近15\%。

NVIDIA也以20亿美元投资入股Nebius，进一步巩固了双方的战略绑定。
Nebius的优势在于其\textbf{欧洲地理位置}——在数据主权日益敏感的监管
环境下，一个总部位于阿姆斯特丹、数据中心分布在芬兰和欧洲各国的
AI云提供商，对需要GDPR合规的欧洲客户具有天然吸引力。从Yandex带来的
工程团队在分布式训练和推理优化方面有着多年积累，这也使得Nebius在技术
深度上区别于那些单纯"买卡卖时间"的竞争对手。

% ===========================================================
\subsection{Together AI：学术殿堂走出的推理引擎}
\label{subsec:together}
% ===========================================================

Together AI的创始团队阵容堪称AI基础设施领域的"全明星"：CEO Vipul
Ved Prakash是连续创业者，联合创始人包括斯坦福教授Christopher
R\'{e}（数据基础设施先驱）和FlashAttention的发明者Tri Dao。这种
学术-工程的深度融合让Together AI从一开始就在推理效率上拥有技术优势。

公司成立于2022年，在2025年完成了3.05亿美元B轮融资，估值33亿美元。
到2026年3月，Together AI的年化营收已突破10亿美元（较2025年中增长
3倍以上），正在洽谈以75亿美元估值融资10亿美元（The Information,
2026年3月）。平台已服务超过45万开发者。

Together AI的定位介于纯GPU云和模型API之间——它不仅提供裸GPU，更提供
\textbf{高度优化的推理和微调服务}。凭借FlashAttention等底层优化技术，
Together AI声称其推理性能在同等硬件上可比竞争对手高出数倍。这种
"基础设施+优化层"的叠加，使其在开源模型部署领域建立了强大的护城河。
一个值得注意的信号是：当越来越多的企业选择部署开源模型（如Llama、
Mistral、DeepSeek）而非调用闭源API时，Together AI这样的"开源友好型"
推理平台将直接受益于这一结构性趋势。

% ===========================================================
\subsection{RunPod：资本效率的极致样本}
\label{subsec:runpod}
% ===========================================================

在一个充斥着数十亿美元融资的赛道中，RunPod是一个令人惊叹的异类。
这家2022年成立于新泽西的公司，\textbf{仅融资2200万美元}，却在2026年
1月宣布年化经常性收入（ARR）突破1.2亿美元，服务超过50万名开发者
（RunPod, 2026年1月）。

RunPod的起源颇具草根色彩——创始团队最初在Reddit的机器学习社区中建立
了用户基础。产品定位直指\textbf{个人开发者和小型团队}：无需预约、
即时可用的GPU实例，价格通常比AWS便宜50--90\%，界面极其简洁。这种
"开发者优先"的策略让RunPod在社区口碑中获得了远超其规模的影响力。

RunPod在Ramp（企业支出管理平台）2026年3月的趋势榜单中被评为
AI代理基础设施领域的顶级增长厂商——这意味着不仅个人开发者，越来越多
的企业客户也在采用其服务。1.2亿美元ARR除以2200万美元总融资，意味着
每融资1美元产生5.45美元年收入——这在GPU云赛道中是无与伦比的资本效率。

RunPod的故事提出了一个尖锐的问题：在Neocloud赛道中，融资规模到底
是优势还是负担？CoreWeave融了数十亿美元，但也背上了数十亿美元的
债务和折旧压力；RunPod轻装上阵，每一块钱都用在了刀刃上。当然，
RunPod的轻模式也意味着它无法承接超大规模训练任务——\textbf{不同的
资本结构决定了不同的客户群和天花板}。

% ===========================================================
\subsection{Vast.ai：GPU市场的"自由贸易区"}
\label{subsec:vastai}
% ===========================================================

如果说RunPod是资本效率的冠军，Vast.ai则在更极端的方向上走得更远。
这家2016年成立于洛杉矶的公司\textbf{仅融资400万美元}，却建立了一个
真正的GPU P2P市场：任何拥有闲置GPU的数据中心或个人都可以将算力挂到
平台上出售，价格由\textbf{市场竞价机制}动态决定。

这种"GPU版Airbnb"的模式意味着Vast.ai几乎不需要自己购买和运维硬件
——它只提供市场基础设施和信任机制。2026年3月，Vast.ai被Ramp和Brex
两大企业支出平台同时评为增长最快的厂商之一（Vast.ai, 2026年3月），
表明其市场模式正在从个人开发者向企业客户渗透。

Vast.ai的价格透明度是其最大卖点：用户可以实时看到不同型号GPU在不同
地区的现货价格，根据成本/性能比自主选择。对于不需要保证SLA的研究和
实验任务，这种弹性模式极具吸引力。值得注意的是，Vast.ai与VAST Data
（AI存储公司，2026年估值\$30B）是完全不同的公司，尽管名字容易混淆。

% ===========================================================
\subsection{Modal：一行代码的Serverless GPU}
\label{subsec:modal}
% ===========================================================

Modal Labs代表了GPU云赛道中最"开发者友好"的一极。创始人Erik
Bernhardsson曾在Spotify担任机器学习负责人，亲历了在传统云上部署
ML模型的痛苦——配置实例、管理容器、处理自动扩缩容，工程开销常常超过
模型本身。2021年，他在纽约创立Modal，目标很简单：\textbf{让GPU计算
像调用函数一样简单}。

Modal的产品体验确实做到了这一点。开发者只需在Python函数上加一个装饰器
（decorator），就能将计算任务自动调度到GPU集群上执行，按实际使用时间
计费，空闲时不产生任何费用。这种Serverless GPU的抽象极大地降低了
使用门槛——它把云GPU从一个"基础设施问题"变成了一个"函数调用问题"。

2025年9月，Modal完成了8700万美元B轮融资，估值11亿美元，正式跻身
独角兽行列。2026年2月，TechCrunch报道Modal正在洽谈以25亿美元估值
融资新一轮（TechCrunch, 2026年2月）。估值在不到半年内翻倍，反映了
市场对Serverless GPU这一范式的强烈看好。

Modal的战略意义在于它预示了GPU云的\textbf{下一阶段}：当底层算力逐渐
商品化，价值将向更高层的抽象和开发者体验转移。正如AWS S3让开发者不再
关心磁盘管理，Modal让开发者不再关心GPU编排。对于正在快速增长的
AI Agent工作负载——需要动态、突发性的GPU推理——Serverless模式有着
天然的架构优势。

% ===========================================================
\subsection{其他重要玩家}
\label{subsec:gpu-cloud-others}
% ===========================================================

\subsubsection{Voltage Park $\to$ Lightning AI}
Voltage Park由Stellar/XRP联合创始人Jed McCaleb的家族基金出资5亿美元
于2023年创立，以"使命驱动"为理念，与NSF等公共机构合作提供低成本学术
算力。2026年1月，Voltage Park与AI开发平台Lightning AI（由PyTorch
Lightning创始人William Falcon领导）完成合并，创建了"第一个为AI而生
的云"。合并后的公司宣称从2024年的1800万美元营收增长至超过5000万美元
（BusinessWire, 2026年1月）。这一组合将GPU硬件资产与开发者工具整合
在一起，试图提供从实验到生产的完整闭环——硬件公司和软件公司各取所需，
是Neocloud整合浪潮的早期范例。

\subsubsection{FluidStack}
FluidStack成立于2017年伦敦，走的是分布式GPU聚合的路线。公司的关键
突破在于2025年6月与Macquarie集团签署的100亿美元债务融资协议，用于在
欧洲部署GPU基础设施。更具战略意义的是，FluidStack成为\textbf{Google
TPU的首个第三方分销伙伴}——在NVIDIA主导的Neocloud生态中，这是一个
罕见的差异化选择。2026年2月，Google据报道正在考虑向FluidStack投资
1亿美元（SiliconANGLE, 2026年2月）。此外，FluidStack参与了法国
政府100亿欧元AI超级计算项目的竞标，试图在欧洲主权AI基础设施领域
占据一席之地。然而，2026年3月也传出FluidStack面临资本压力的消息
（AInvest, 2026年3月），说明在重资产赛道中，即便有百亿美元级的
债务承诺，现金流管理依然是生死攸关的挑战。

\subsubsection{SF Compute}
SF Compute成立于2023年旧金山，定位是\textbf{GPU算力交易所}——一个
标准化的商品市场，买家和卖家可以像交易大宗商品一样交易GPU计算时间。
2025年12月，公司完成4000万美元A轮融资，估值3亿美元（DCD, 2025年
12月）。其模式与Vast.ai的P2P市场有相似之处，但更强调机构级的合约
标准化和价格发现机制，瞄准的是大型企业和GPU持有者之间的批发交易。
值得关注的还有Compute Exchange，这是一个2026年出现的新玩家，试图
建立更接近金融交易所的GPU期货/现货市场。

\subsubsection{新进入者：Andromeda AI 与 Hosted.ai}
2026年3月，GPU云赛道仍在吸引新的融资。Andromeda AI以15亿美元估值
完成了由Paradigm领投的6000万美元融资，聚焦按需弹性GPU租赁
（SiliconANGLE, 2026年3月）。同月，Hosted.ai完成了1900万美元种子轮，
其软件平台旨在帮助现有数据中心运营商更高效地池化和调度GPU资源——
不是建新的数据中心，而是让已有的数据中心跑得更好。这些新进入者表明，
尽管头部玩家已经形成了显著的规模优势，Neocloud赛道的\textbf{垂直
细分机会}仍然在不断涌现。

% ===========================================================
% Comparison table
% ===========================================================

\begin{table}[htbp]
\centering
\caption{GPU 云创业公司全景对比（截至2026年3月）}
\label{tab:gpu-cloud-comparison}
\small
\begin{tabularx}{\textwidth}{l c r r X}
\toprule
\rowcolor{figLightGray}
\textbf{公司} & \textbf{成立} & \textbf{总融资/市值}
  & \textbf{2025营收} & \textbf{独特角度} \\
\midrule
CoreWeave   & 2017 & 市值\$43.5B
  & \$5.13B & 加密转型; NVIDIA Exemplar; 收购W\&B \\
Lambda Labs & 2012 & 总融\$2.3B
  & $\sim$\$505M & 学术根基; ``超级智能云'' \\
Crusoe      & 2018 & 估值\$10B+
  & 未披露 & 废气发电; Stargate联盟 \\
Nebius      & 2024 & 市值$\sim$\$23B
  & \$529.8M & Yandex分拆; 欧洲GDPR; Meta \$27B \\
Together AI & 2022 & \$3.3B$\to$\$7.5B
  & ARR$\sim$\$1B & FlashAttention; 推理优化 \\
Modal       & 2021 & \$1.1B$\to$\$2.5B
  & 未披露 & Serverless GPU; 极致DX \\
RunPod      & 2022 & 仅融\$22M
  & ARR \$120M & 极致资本效率; 社区驱动 \\
Vast.ai     & 2016 & 仅融\$4M
  & 未披露 & P2P市场; 竞价定价 \\
Lightning/VP& 2023 & 基金\$500M
  & $\sim$\$50M & 使命驱动; 硬件+软件合并 \\
FluidStack  & 2017 & 债务\$10B
  & 未披露 & Google TPU分销; 欧洲布局 \\
SF Compute  & 2023 & 估值\$300M
  & 未披露 & 算力交易所; 机构化 \\
\bottomrule
\end{tabularx}
\end{table}

% ===========================================================
% Market structure analysis
% ===========================================================

\subsection{赛道结构与竞争逻辑}
\label{subsec:gpu-cloud-landscape}

纵观上述11家公司（及新进入者），GPU云赛道已形成清晰的三层结构：

\textbf{第一层：资本密集型巨头}。CoreWeave和Nebius代表了"自建数据
中心+长期合约"的重资产模式。这两家公司的共同特征是：融资规模以数十亿
美元计，客户集中于少数超大型AI实验室，营收增速惊人但资本支出同样惊人。
CoreWeave的668亿美元积压订单看似令人安心，但其中大部分来自少数几个
客户，违约风险不可忽视。

\textbf{第二层：差异化中型玩家}。Lambda、Crusoe、Together AI和
FluidStack各自找到了独特的定位——学术基因、绿色能源、推理优化、
欧洲市场。这一层的竞争关键在于：\textbf{你的差异化是否足够深，以至于
客户不会在H100价格下降时简单地切换到更便宜的替代方案？}

\textbf{第三层：轻资产创新者}。RunPod、Vast.ai、Modal和SF Compute
则代表了另一种哲学——不拥有硬件（或少量拥有），而是在市场机制、
开发者体验或金融抽象层面创新。这些公司的资本效率远高于前两层，但
面临的问题是\textbf{规模天花板}和\textbf{企业客户的信任度}。

\begin{corebox}[GPU 云的价值链分层]
从底层到上层，GPU云赛道的价值链可以这样理解：
\begin{enumerate}
  \item \textbf{物理层}：数据中心、电力、冷却（Crusoe、Voltage Park）
  \item \textbf{硬件层}：GPU/TPU采购与集群化
    （CoreWeave、Lambda、Nebius）
  \item \textbf{市场层}：算力的撮合与定价
    （Vast.ai、SF Compute）
  \item \textbf{编排层}：Serverless抽象与自动扩缩
    （Modal、RunPod）
  \item \textbf{优化层}：推理/训练性能优化
    （Together AI、FluidStack）
  \item \textbf{平台层}：MLOps工具与开发者体验
    （W\&B/CoreWeave、Lightning AI）
\end{enumerate}
每家公司都在试图向相邻层扩展——CoreWeave收购W\&B是从硬件层向平台层
扩展，Lightning AI合并Voltage Park是从编排层向物理层扩展。这种
\textbf{垂直整合}的趋势将在未来两年加速。
\end{corebox}

\begin{warnbox}[风险警告：GPU 云赛道的三大结构性风险]
\begin{enumerate}
  \item \textbf{NVIDIA 依赖症}。几乎所有Neocloud的核心资产都是
    NVIDIA GPU。NVIDIA通过Exemplar Cloud、DGX Cloud等项目对算力分配
    拥有巨大影响力。当NVIDIA自己开始直接向终端客户销售云服务时，
    Neocloud是合作伙伴还是竞争对手？FluidStack选择Google TPU是对
    这一风险的有意对冲。
  \item \textbf{资本消耗的无底洞}。一块H100的采购成本约
    \$25,000--30,000，一个万卡集群的前期投资就超过3亿美元，还不算
    电力、冷却和运维。这意味着GPU云公司需要持续进行大规模融资或
    举债——当利率上升或AI投资情绪降温时，现金流压力将急剧放大。
    FluidStack 2026年3月的资本压力即为早期警示信号。
  \item \textbf{客户集中度}。CoreWeave的前五大客户贡献了绝大部分
    营收；Nebius的270亿美元合同来自单一客户Meta。在超大规模客户
    拥有极强议价能力的市场中，GPU云提供商本质上是在\textbf{用规模
    换利润率}。
\end{enumerate}
\end{warnbox}

这个赛道最令人着迷的张力在于：GPU算力正在经历\textbf{从稀缺资源到
基础设施商品}的转变。2023年，一块H100的租赁价格可达每小时\$8--12；
到2026年初，价格已降至\$2--4/小时。随着NVIDIA B系列GPU和AMD MI300X
的大规模出货，供给端的瓶颈正在缓解。

这对不同层级的玩家意味着截然不同的未来。对CoreWeave和Nebius来说，
长期锁定合约提供了缓冲垫，但终究无法逃避\textbf{算力贬值}的趋势。
对RunPod和Vast.ai这样的市场模式来说，价格下降反而可能扩大需求——当
GPU足够便宜时，长尾开发者和中小企业将加速上云。对Modal这样的抽象层
玩家来说，底层硬件的商品化是最大的利好——正如云存储的商品化让S3成为
基础设施的一部分，GPU的商品化将让Serverless GPU成为默认的计算范式。

未来两三年内，这个赛道几乎必然会经历一轮整合。Lightning AI与
Voltage Park的合并已经拉开了序幕。那些无法在规模、差异化或资本效率上
建立护城河的玩家，将面临被收购或淘汰的命运。而最终胜出的公司，很可能
不是拥有最多GPU的那个，而是\textbf{最懂得如何让开发者忘记GPU存在}
的那个。

% ===========================================================
\subsection{长尾玩家：小众与独立 GPU 云}
\label{subsec:gpu-cloud-longtail}
% ===========================================================

在头部和中型 GPU 云之外，一批更小、更聚焦的玩家正在覆盖主流服务商忽略
的缝隙市场。它们的共同特征是：融资规模不大，但在特定地理区域、特定用户
群或特定技术路线上建立了独特的存在感。

\subsubsection{JarvisLabs.ai}
JarvisLabs（印度）是面向个人开发者和学生的超低成本 GPU 云，A100 实例
低至 \$1.29/hr，H100 低至 \$2.99/hr，按分钟计费，90 秒即可启动。
平台聚焦深度学习实验和 Stable Diffusion 工作负载，
截至 2026 年初已服务超过 27,000 名 AI 开发者。
JarvisLabs 以极致的简洁性和低进入门槛吸引了大量 fast.ai 学员和
独立研究者——对于只需要几小时 GPU 时间跑一次实验的用户来说，
它比 AWS 和 RunPod 都更轻便。

\subsubsection{TensorDock}
TensorDock（2019 年，美国）定位为``极致性价比 GPU 云''，
H100 低至 \$2.25/hr，RTX 4090 低至 \$0.35/hr。
2025 年被 Voltage Park 收购后纳入 Lightning AI 生态，
但品牌和平台保持独立运营。TensorDock 的核心优势是
\textbf{全球分布式 GPU 舰队}——覆盖多个 Tier 3/4 数据中心，
并提供业界最广的 GPU 型号选择。无需预约、无最低消费、
最低充值仅 \$5 的模式让它在独立开发者中极具吸引力。

\subsubsection{Hyperbolic}
Hyperbolic（2024 年，Berkeley）自称``开放访问 AI 云''，
由 Databricks 联合创始人 Reynold Xin 担任顾问。
商业模式结合了\textbf{推理 API + GPU 租赁 + 算力供给}三位一体：
用户既可以调用推理 API，也可以租用裸 GPU，
还可以将自己的闲置 GPU 接入平台赚取收益。
这种三面市场的模式试图同时解决供给侧和需求侧的匹配问题，
定价号称比主流云商低 80\% 以上。

\subsubsection{Salad（SaladCloud）}
Salad（2018 年，Salt Lake City）是 GPU 云赛道中最另类的存在：
它将全球\textbf{数十万台游戏 PC 的闲置 GPU} 聚合为一个分布式计算网络。
玩家在不打游戏时运行 Salad 客户端，将自己的 RTX 显卡算力出租给
AI 推理任务，按贡献获得报酬。公司累计融资 \$27.5M（Series A），
2025 年营收据报接近 \$200M 水平。
这种``消费级 GPU 众包''模式在数据中心 GPU 供不应求的时期提供了
一个非传统的补充方案，但可靠性和 SLA 保证是其天然劣势。

\subsubsection{Prime Intellect}
Prime Intellect（San Francisco）聚焦\textbf{去中心化 AI 训练}——
将全球分散的小型 GPU 集群通过自研的分布式训练协议组织起来，
使开源社区无需百万美元预算也能训练大模型。
其技术栈针对跨数据中心、高延迟网络环境下的分布式训练做了深度优化，
是 DiLoCo（Distributed Low-Communication）等新型分布式算法的
重要实践者。Prime Intellect 代表了 GPU 云的另一种可能：
\textbf{不是建数据中心，而是把已有的算力拼起来}。

\subsubsection{Shadeform}
Shadeform（2023 年，San Francisco，Y Combinator S23）是 GPU 云的
``聚合搜索引擎''：通过统一 API 接入 20 多家 GPU 云供应商
（包括 Lambda、Nebius、Crusoe 等），让用户一键比价并跨云部署实例。
团队仅 5--7 人，融资约 \$500K，但已实现约 \$1M 营收。
Shadeform 不拥有任何 GPU，其价值完全在于\textbf{抽象层}——
类似机票比价网站 Kayak 之于航空公司。如果 GPU 云继续碎片化，
这种聚合器模式将变得更有价值。

\subsubsection{Genesis Cloud}
Genesis Cloud（2018 年，Munich）是欧洲主权 AI 云的代表之一。
累计融资仅 \$6.6M，但凭借北欧绿色数据中心和 NVIDIA 参考架构，
成为欧盟 AI Champions Initiative 的成员，服务于
对 GDPR 合规和数据主权有严格要求的欧洲企业客户。
定价较传统云厂商低约 80\%，H100 集群已部署就绪。

\subsubsection{E2E Networks}
E2E Networks（印度，NSE 上市：E2ENET）是印度本土最大的 AI GPU 云，
获得 IndiaAI Mission 下 ₹177 亿卢比政府合同，为印度首个主权基础
AI 模型提供 GPU 算力。2026 年 1 月，E2E 采购了 1024 块 NVIDIA B200
部署于 Chennai 数据中心。作为印度 AI 基础设施的国家队选手，
E2E 代表了 GPU 云赛道的\textbf{地理多元化}趋势——
当每个国家都在追求``AI 主权''时，本土 GPU 云提供商将获得
政策扶持和数据合规方面的天然优势。

\subsubsection{Oblivus Cloud}
Oblivus（London）是面向欧洲市场的小众 GPU 云，
提供 RTX 5090、H100 等多种型号，覆盖 9 个数据中心、
超过 18,000 块 GPU。号称比主流云商便宜 80\%，
无配额、无验证流程、按需部署。虽然团队极小
（LinkedIn 显示仅数人），但其低门槛策略在预算有限的
欧洲初创公司和独立研究者中找到了用户群。


% §2.2 AI Chip Startups (Western + Chinese)
\section{AI 芯片新势力}
\label{sec:infra-chips}

% ch02-s2a-chips-west.tex  — §2.2 Western AI Chip Startups
% Part of Chapter 2: Infrastructure Layer
% Input from: ch02-infrastructure.tex (or main file)

\subsection{挑战 NVIDIA：西方 AI 芯片创业浪潮}
\label{subsec:chips-west}

2024 年，NVIDIA 数据中心 GPU 营收突破 \$470 亿，
占全球 AI 训练芯片市场份额超过 90\%。
CUDA 生态经过十余年积累，已成为深度学习事实标准——
从 PyTorch 到 TensorRT，从 cuDNN 到 NCCL，
整个软件栈深度绑定 NVIDIA 硬件。
在这个看似铁板一块的市场里，
一个问题自然浮现：\textbf{为什么还有人前赴后继地造芯片？}

答案藏在三个结构性裂缝中。
其一，训练与推理的\textbf{架构需求正在分化}。
训练需要极致的浮点算力和跨节点通信带宽，
而推理——尤其是大规模部署的 LLM 推理——
瓶颈在于内存带宽和延迟确定性
（参见 \textit{Emergence} 第九章对推理瓶颈的分析）。
通用 GPU 为训练优化的架构在推理场景中存在大量浪费。
其二，\textbf{成本压力倒逼替代方案}。
当 AI 公司年花费数十亿美元购买 H100/B200 集群时，
哪怕 10\% 的性能功耗比提升都意味着数亿美元的节省。
其三，\textbf{供应链风险}催生多元化需求。
单一供应商的地缘政治风险（美国对华出口管制）
和产能瓶颈（CoWoS 封装排队）
让超大规模云厂商（Google TPU、AWS Trainium、Microsoft Maia）
和创业公司同时寻求替代路径。

以下按技术路线梳理西方主要 AI 芯片创业公司。

\begin{corebox}[AI 芯片技术路线全景]
\begin{itemize}[noitemsep]
  \item \textbf{晶圆级计算}——将整片晶圆做成一颗芯片，
        突破 reticle 面积限制（Cerebras）
  \item \textbf{确定性推理}——时序可预测的数据流架构，
        消除 GPU 调度不确定性（Groq）
  \item \textbf{可重构数据流}——编译器驱动的空间架构，
        针对大模型数据流优化（SambaNova）
  \item \textbf{RISC-V + AI}——开源指令集与 AI 加速融合，
        芯片设计 + IP 授权双线策略（Tenstorrent）
  \item \textbf{存内计算}——将计算搬到存储旁边，
        突破``内存墙''瓶颈（d-Matrix）
  \item \textbf{Transformer 专用 ASIC}——
        烧死一切非 Transformer 架构灵活性，
        押注单一模型族（Etched、MatX）
  \item \textbf{边缘 AI 处理器}——
        低功耗、低延迟的端侧推理芯片（Hailo）
\end{itemize}
\end{corebox}

\subsubsection{Cerebras：一整片晶圆就是一颗芯片}

Andrew Feldman 在 2015 年创立 Cerebras 时，
提出了一个大胆到近乎荒谬的想法：
\textbf{为什么要把晶圆切成几百颗小芯片？直接用整片晶圆做计算。}
传统芯片受限于光刻机的 reticle 尺寸（约 $26 \times 33$\,mm），
单颗芯片面积不超过 $\sim$850\,mm$^2$。
Cerebras 的 Wafer-Scale Engine（WSE）则占据整片 300mm 晶圆，
面积达 46,225\,mm$^2$——是 NVIDIA B200 的 50 倍以上。

第三代 WSE-3 采用 TSMC 5nm 工艺，
集成 4 万亿晶体管、90 万个 AI 优化核心，
片上 SRAM 达 44\,GB。
关键在于，WSE 将\textbf{计算和存储放在同一片硅上}，
省去了传统 GPU 集群中最昂贵的跨芯片通信开销。
对于稀疏模型和大规模科学计算，
这种架构的效率优势尤为显著
（参见 \textit{Emergence} 第五章对训练通信瓶颈的讨论）。

商业化方面，Cerebras 走了一条独特路径：
与 AWS 合作提供云端推理服务，
同时面向国家实验室和主权 AI 项目销售整机系统。
2023 年营收 \$78.7M，虽然远不及 NVIDIA，
但在 AI 芯片创业公司中已属领先。
2024 年底，Cerebras 首次递交 IPO 申请，
但因中东客户（Group 42）的地缘政治审查而搁置。
2026 年 2 月，公司完成 Series H 融资 \$10 亿，
估值达 \$231 亿，随后聘请 Morgan Stanley 牵头，
\textbf{目标 2026 年 4 月重启 IPO，计划融资 \$20 亿}。
若成功上市，Cerebras 将成为 AI 芯片创业公司中的首个 IPO 标杆。

\subsubsection{Groq：从独立创业到 NVIDIA 旗下}

Jonathan Ross 是 Google TPU 的发明者之一，
2016 年离开 Google 创立 Groq，
带着一个清晰的技术判断：
\textbf{GPU 的调度不确定性是推理的天敌}。
GPU 的 warp scheduler 和缓存层级虽然灵活，
但在实时推理中引入了不可预测的延迟抖动。
Groq 的 Language Processing Unit（LPU）采用
\textbf{时序确定性（time-deterministic）}数据流架构：
每条指令在编译时就确定了执行时间，
运行时不存在缓存未命中或调度冲突。

这种确定性带来了惊人的推理速度：
Groq 的 LPU 系统在 Llama 2 70B 上实现了超过 500 tok/s 的单用户吞吐，
延迟几乎为零抖动。
GroqCloud API 上线后迅速吸引了超过 250 万开发者，
``Groq fast'' 一度成为 AI 社区的流行语。

然而，Groq 的故事在 2025 年 12 月迎来了戏剧性转折：
\textbf{NVIDIA 以 \$200 亿完成对 Groq 的收购}。
这并非一次简单的消灭竞争对手——
在 2026 年 3 月的 GTC 大会上，Jensen Huang 亲自解释了战略逻辑：
LPU 不是取代 GPU，而是\textbf{补全}GPU。
NVIDIA 随即发布了 LPX 机架系统，
将 Groq 的 LPU 与自家 Vera Rubin GPU 协同部署：
GPU 处理训练和 prefill，LPU 负责 decode 推理。
更值得关注的是，NVIDIA 正在为中国市场准备合规版 Groq 芯片，
将确定性推理架构的影响力进一步扩展。

Groq 的被收购验证了一个产业规律：
在 AI 芯片领域，技术差异化可以创造真实价值，
但\textbf{独立生存}需要的不仅是好架构——
还需要与 CUDA 抗衡的软件生态。

\subsubsection{Tenstorrent：Jim Keller 的开源芯片野心}

如果说 AI 芯片行业有一位``传奇设计师''，
那就是 Jim Keller——
AMD Zen 架构、Apple A4/A5 处理器、
Tesla 自动驾驶芯片 FSD 的核心设计者。
2020 年，Keller 加入多伦多创业公司 Tenstorrent 担任 CEO，
带来了一个与众不同的战略：
\textbf{不仅造芯片，还做开源 AI 计算平台}。

Tenstorrent 的独特之处在于双模架构：
RISC-V 通用处理器核心 + Tensix AI 加速核心共存于同一芯片。
这意味着芯片既可以运行操作系统和通用代码，
也可以执行高效的 AI 推理——
而整个软件栈基于开源 RISC-V ISA 和开源编译器。
这直接对标 NVIDIA 的 CUDA 锁定策略，
试图为 AI 计算建立一个``Linux 式''的开放生态。

2025 年底，Tenstorrent 完成 Series D 融资 \$6.93 亿，
估值达 \$20 亿。
公司同时推行\textbf{芯片 + IP 授权}双线策略：
一方面出售自研的 Wormhole 和 Blackhole 系列芯片，
另一方面将 RISC-V + AI 的 IP 模块授权给第三方厂商。
2026 年 3 月，Tenstorrent 发布了 TT-QuietBox 2（Blackhole），
这是首款基于 RISC-V 的液冷 AI 工作站，
能在桌面端本地运行 1200 亿参数的模型，
全部基于开源软件栈，起价 \$9,999。

Keller 的赌注清晰：
\textbf{AI 芯片的终局不是谁的硬件更快，
而是谁的生态更开放}。
如果 RISC-V 在 AI 领域复制 Linux 在服务器市场的成功，
Tenstorrent 将成为这个生态的 Red Hat。

\subsubsection{SambaNova：数据流架构的坚守者}

SambaNova 由斯坦福大学教授 Kunle Olukotun 和 Chris R\'{e}
（Flash\-Attention 的学术推动者之一）
联合 Rodrigo Liang 于 2017 年创立，
是学术创业在 AI 芯片领域的典型代表。
其 Reconfigurable Dataflow Unit（RDU）采用编译器驱动的空间架构，
在编译阶段将模型计算图直接映射到硬件数据流管线上，
实现了接近理论峰值的硬件利用率。

2026 年 2 月，SambaNova 发布了 SN50 推理芯片，
号称在 agentic AI 场景下比 NVIDIA B200 快 5 倍、能效高 8 倍，
并宣布与 Intel（代工合作）和 SoftBank（部署合作）达成战略伙伴关系。
同期完成 Series E 融资 \$3.5 亿，累计融资约 \$14.8 亿。

然而，SambaNova 的估值轨迹揭示了 AI 芯片创业的残酷现实：
2021 年巅峰期估值达 \$51 亿，
2026 年的 Series E 估值却回落至约 \$21 亿——
\textbf{缩水超过 58\%}。
这并非技术失败，而是市场给出的信号：
在 NVIDIA 主导的格局下，
即使拥有顶尖的架构创新，
商业化爬坡的难度远超学术论文的预期。

\subsubsection{d-Matrix：在存储旁边做计算}

d-Matrix（2019 年成立于 Santa Clara）瞄准了 AI 推理的核心瓶颈：
\textbf{内存墙}。
正如 \textit{Emergence} 第九章所分析的，
LLM decode 阶段 99\% 以上的时间花在等待数据从显存搬运到计算单元。
d-Matrix 的解决方案是\textbf{数字存内计算}（Digital In-Memory Compute, DIMC）：
将计算逻辑直接嵌入 SRAM 阵列中，
让数据``就地''完成矩阵乘法，
从根本上消除数据搬运开销。

其 Corsair 平台基于 chiplet 架构，
将多个 DIMC 计算核心通过高速互联组合。
2024 年底完成 Series C 融资 \$2.75 亿，估值达 \$20 亿。
2025--2026 年间，d-Matrix 进一步推出 3D DRAM 堆叠方案（3DIMC），
与 Alchip 合作开发全球首个 3D DRAM 存内计算芯片，
目标是将存内计算的容量从 SRAM 级别扩展到 DRAM 级别，
以适配百亿甚至万亿参数模型的推理需求。

\subsubsection{Etched 与 MatX：押注 Transformer 的极端赌徒}

如果说上述公司还在``通用''与``专用''之间寻找平衡，
Etched 和 MatX 则走到了专用化的极端：
\textbf{只为 Transformer 而生，别的架构一律不支持。}

\textbf{Etched}（2022 年，两位 Harvard 辍学生创立于 San Jose）
开发的 Sohu 芯片是一颗纯粹的 Transformer ASIC。
它将 Attention 计算、KV Cache 管理等
Transformer 特有操作直接烧入硅片，
号称比 NVIDIA H100 快 10--20 倍。
Sohu 采用 reticle-limit 级芯片面积，配备 144\,GB HBM3E。
2026 年 1 月，Etched 完成 \$5 亿融资（Peter Thiel 参投），
估值达 \$50 亿——
但截至 2026 年 3 月，\textbf{芯片尚未出货，
没有任何独立基准测试}。

\textbf{MatX}（2022 年，ex-Google TPU 团队创立于 Mountain View）
走的是类似路线：为 LLM 推理设计专用芯片 MatX One，
核心创新是``可拆分 systolic array''——
既保持大规模矩阵运算的面积与能效优势，
又能灵活适配不同形状的小矩阵。
2026 年 2 月完成 Series B \$5 亿融资，
由 Jane Street 和 Leopold Aschenbrenner 的
Situational Awareness 基金领投。
MatX 计划 2026 年完成芯片设计、
2027 年通过 TSMC 量产出货。

两家公司的共同赌注是：
Transformer 架构将在未来 5--10 年继续主导 AI，
因此牺牲通用性换取 10 倍以上的专用加速是值得的。
但风险同样巨大——
如果 State Space Model（如 Mamba）或其他新架构崛起，
这些芯片将变成昂贵的废硅。

\subsubsection{Hailo：边缘 AI 的隐形冠军}

并非所有 AI 芯片都瞄准数据中心。
以色列公司 Hailo（2017 年，Tel Aviv）专注于\textbf{边缘 AI 推理}：
自动驾驶、安防摄像、工业检测等需要低功耗实时推理的场景。
累计融资 \$3.4 亿，估值约 \$12 亿。

Hailo 的芯片架构核心是``Structure-Defined Dataflow''——
一种面向神经网络拓扑自动优化数据路径的空间架构，
在 10W 以下功耗实现 26\,TOPS 的算力。
2026 年 1 月，Hailo 发布 Hailo-10H 处理器，
支持边缘端的\textbf{生成式 AI 推理}（LLM/VLM），
并与 Raspberry Pi 合作推出 AI HAT+ 2 模组，
首次将生成式 AI 能力带入 \$70 价位的开发板。

在数据中心巨头混战的喧嚣之外，
Hailo 代表了 AI 芯片另一条可行的创业路径：
\textbf{避开 NVIDIA 的主战场，在边缘场景建立垂直壁垒}。

\begin{warnbox}[高估值陷阱：Graphcore 的教训]
在讨论 AI 芯片创业的光明前景时，
不能回避一个沉痛的案例：\textbf{Graphcore}。

这家 2016 年成立于英国 Bristol 的公司
曾是欧洲 AI 芯片的旗帜，
其 Intelligence Processing Unit（IPU）采用独特的``Bulk Synchronous Parallel''
执行模型，针对图神经网络和稀疏计算做了深度优化。
巅峰时期估值达 \$28 亿，累计融资 \$7.67 亿，
投资者包括微软、三星、红杉等顶级机构。

然而，当 LLM 浪潮席卷一切时，
IPU 架构在 Transformer 密集计算上的优势不再明显，
客户增长停滞，现金流告急。
2024 年 7 月，SoftBank 以约 \$4 亿收购 Graphcore——
\textbf{较峰值估值暴跌 86\%}，
投资者损失惨重。

Graphcore 的教训直指 AI 芯片创业的核心风险：
\textbf{技术押注与架构范式的错配}。
IPU 为前 LLM 时代的多样化模型架构而设计，
却在 Transformer 一统天下的新格局中失去了差异化价值。
今天押注 Transformer ASIC 的 Etched 和 MatX，
是否会在下一次架构变革中重蹈覆辙？
这个问题值得每一位投资者深思。
\end{warnbox}

\subsubsection{西方 AI 芯片创业全景对比}

表~\ref{tab:west-chips-comparison} 汇总了主要西方 AI 芯片创业公司的关键指标。

\begin{table}[htbp]
\centering
\caption{西方 AI 芯片创业公司对比（截至 2026 年 3 月）}
\label{tab:west-chips-comparison}
\small
\begin{tabularx}{\textwidth}{@{}l l l l r r l@{}}
\toprule
\rowcolor{figLightGray}
\textbf{公司} & \textbf{成立} & \textbf{核心架构} & \textbf{目标场景} &
  \textbf{总融资} & \textbf{估值} & \textbf{状态} \\
\midrule
Cerebras     & 2015 & 晶圆级 WSE         & 训练 + 推理   & \$1B+  & \$23.1B & 待 IPO \\
Groq         & 2016 & 确定性 LPU         & 推理          & ---    & \$20B 收购 & NVIDIA 旗下 \\
Tenstorrent  & 2016 & RISC-V + Tensix    & 推理 + 边缘   & \$693M & \$2B    & 产品出货 \\
SambaNova    & 2017 & 数据流 RDU         & 训练 + 推理   & \$1.48B & $\sim$\$2.1B & 估值缩水 \\
d-Matrix     & 2019 & 存内计算 DIMC      & 推理          & \$275M & \$2B    & 研发中 \\
Etched       & 2022 & Transformer ASIC   & 推理          & \$625M & \$5B    & 未出货 \\
MatX         & 2022 & LLM 专用 systolic  & 推理          & \$500M+ & ---    & 未出货 \\
Hailo        & 2017 & 边缘数据流         & 边缘推理      & \$340M & \$1.2B  & 产品出货 \\
Graphcore    & 2016 & IPU (BSP)          & 训练          & \$767M & $\sim$\$0.4B 收购 & SoftBank 旗下 \\
\bottomrule
\end{tabularx}
\end{table}

从表中可以看出几个清晰的趋势：
\textbf{第一}，推理正在取代训练成为芯片创业的主赛道——
九家公司中有七家将推理作为核心场景或主要场景，
这与 LLM 部署成本日益超过训练成本的产业趋势一致。
\textbf{第二}，收购整合正在加速——
Groq 被 NVIDIA 以 \$200 亿收购，Graphcore 被 SoftBank 低价收入，
这说明 AI 芯片市场正在从``百花齐放''走向``赢者通吃''与``被收编''的分化格局。
\textbf{第三}，高估值与未出货之间的张力令人警惕——
Etched（\$50 亿）和 MatX 都在芯片未量产的情况下获得巨额融资，
这既反映了资本对 NVIDIA 替代方案的渴望，
也埋下了 Graphcore 式泡沫破裂的隐患。

对于 AI 创业生态而言，芯片层的竞争格局直接决定了上层的可能性：
更多元的芯片供给意味着更低的推理成本，
更低的推理成本意味着更多的 AI 应用可以在经济上成立。
这正是为什么在 NVIDIA 看似不可战胜的今天，
资本仍然源源不断地涌入 AI 芯片创业——
\textbf{它们押注的不是打败 NVIDIA，
而是把 AI 计算这块蛋糕做得足够大，
大到连 10\% 的份额也值得万亿美元}。

% ===========================================================
\subsubsection{更多西方与亚洲 AI 芯片创业公司}
% ===========================================================

上述公司之外，还有一批值得关注的 AI 芯片创业力量——
它们或聚焦边缘场景，或来自韩国的半导体创业浪潮，
或在特定技术缝隙中寻找生存空间。

\paragraph{韩国``K-NVIDIA''军团。}
韩国政府在 2026 年宣布每年投入 10 万亿韩元（约 \$75 亿）
扶持本土 AI 芯片产业，并推出 K-Perf 统一性能基准。
三家核心创业公司正在崛起：

\textbf{Rebellions}（2020 年，首尔）
在 2024 年 12 月与 SK Telecom 旗下的 SAPEON Korea 完成合并，
成为\textbf{韩国首个 AI 芯片独角兽}。
合并后累计融资约 \$3.49 亿（Series C，2025 年 10 月），
员工增至 167 人，年营收约 \$7300 万。
其 REBEL-Quad 是一颗基于 4nm 的 peta-scale AI ASIC，
采用数据流架构 + HBM3E + chiplet 设计，
面向数据中心推理，已签下沙特、日本、法国等海外客户。
Rebellions 正积极考虑赴美 IPO。

\textbf{FuriosaAI}（首尔）
专注数据中心 AI 推理，旗舰芯片 RNGD（读作``Renegade''）
采用 Tensor Contraction Processor 架构，
在 LG AI Research 的 EXAONE 平台上实现了
比 GPU 高 2.25 倍的推理能效。
2025 年 7 月完成 \$1.25 亿 Series C bridge 融资
（累计 \$2.46 亿，估值 \$7.35 亿），
2026 年 1 月据报道正在融资 \$3--5 亿 Series D。
FuriosaAI 是韩国 AI 芯片创业中
\textbf{最接近规模化商业出货}的公司。

\textbf{DEEPX}（2018 年，韩国）
聚焦\textbf{端侧/边缘 AI 芯片}（NPU），
提供四款商用芯片产品线，覆盖机器人、智慧城市、安防等场景。
累计融资 \$1.04 亿（Series C），员工 65 人，营收约 \$840 万。
2025 年 DEEPX 在中国和东南亚签下重要订单，
成为韩国 AI 芯片跨境扩张的先行者。

\paragraph{Blaize：边缘 AI 芯片的 SPAC 上市先驱。}
Blaize（2011 年，El Dorado Hills，前身 ThinCI）是
\textbf{第一家通过 SPAC 在纳斯达克上市的 AI 芯片创业公司}
（2025 年 1 月，NASDAQ: BZAI）。
由前 Intel 工程师创立，累计融资 \$3.35 亿，
投资者包括 Samsung 和 Mercedes-Benz。
Blaize 的可编程图计算架构（Graph Streaming Processor）
面向边缘场景——智能零售、国防、工业视觉，
年营收约 \$4670 万。
上市的象征意义大于财务意义：
它证明了 AI 芯片创业公司可以走通资本市场退出路径，
尽管道路漫长（从成立到上市历时 14 年）。

\paragraph{Kneron（耐能）：从圣地亚哥到全球边缘。}
Kneron（2015 年，San Diego）由刘峻诚（Albert Liu）创立，
累计融资 \$2.12 亿，专注\textbf{全栈边缘 AI}——
从 NPU 芯片到编译器到应用 SDK 的垂直整合。
2025 年 11 月发布的 KL1140 芯片是\textbf{首款能在设备端
运行完整 Mamba 网络的 NPU}，
号称比云端方案成本降低 10 倍、能效提升 3 倍。
Kneron 与 Qualcomm、MediaTek 等平台合作，
客户遍及台湾、中国、韩国的消费电子和安防领域——
与 Hailo 类似，Kneron 选择了\textbf{避开数据中心主战场，
在端侧 AI 建立差异化壁垒}。


% ============================================================
% §2.2b  Chinese AI Chip Companies — National Substitution
% ============================================================
% This file is \input from ch02-infrastructure.tex under
% \section{AI Chip Startups} after the Western chip section.

\subsection{国产替代：中国 AI 芯片的突围之路}
\label{subsec:chips-china}

被封锁反而激发了中国 AI 芯片最猛烈的创业潮。2022 年 10 月，美国商务部工业与安全局（BIS）发布对华半导体出口管制新规，限制 NVIDIA A100/H100 等高端 GPU 对华销售；2023 年 10 月管制再度升级，连"合规特供版"H800/A800 也被封堵；2024 年 12 月，BIS 将 140 家中国半导体实体列入清单，并将 HBM 纳入管制范畴；2026 年 3 月，特朗普政府更起草全球性 AI 芯片出口许可制草案，要求英伟达、AMD 向任何国家出口先进 AI 加速器均须逐案审批。每一轮制裁本意在于遏制中国 AI 算力发展，结果却如同压力锅——封得越紧，国产替代的引擎转得越猛。

据 Bernstein Research 2025 年 12 月报告预测，华为昇腾 AI 芯片将于 2026 年占据中国 AI 芯片市场 50\% 份额，而英伟达则将从 2024 年的 39\% 骤降至 8\%。摩根士丹利 2026 年 3 月报告更进一步：中国 AI GPU 自给率将从 2024 年的 34\% 飙升至 2027 年的 82\%，本土 AI 芯片收入上调至 1800 亿元人民币，2024--2027 年复合年增长率（CAGR）高达 62\%。在此背景下，一批中国 AI 芯片公司在 2025--2026 年密集登陆资本市场，形成了前所未有的"集体上市潮"。

\begin{corebox}[中国 AI 芯片格局：一超多强]
\textbf{一超：华为昇腾。} 从 2018 年昇腾 310 发布起步，历经 910/910B/910C 迭代，2026 年规划 950/960/970 三大系列。2026 年市占率预计 50\%，定义了国产 AI 算力的"基准线"。中国移动宣布 2028 年前部署 10 万张国产 GPU，全部采用本土方案。

\textbf{多强阵营}由三条主线构成：
\begin{itemize}
  \item \textbf{云端训推：}寒武纪（思元系列 ASIC）、海光信息（DCU GPU + C86 CPU）、昆仑芯（百度体系）、燧原科技（邃思系列）；
  \item \textbf{GPU 四小龙：}壁仞科技、摩尔线程、沐曦、天数智芯——四家均在 2025 年 12 月至 2026 年 1 月间完成 IPO，是中国半导体史上罕见的"四连上市"；
  \item \textbf{垂直赛道：}地平线（车载 AI 芯片）、算能（TPU + RISC-V）。
\end{itemize}
\end{corebox}

\begin{whybox}[为什么出口管制反而加速了国产替代？]
表面看，管制切断了中国企业获取 NVIDIA H100/H200 的渠道，理应减缓 AI 发展。但实际效果恰好相反，原因有三：

\textbf{第一，需求刚性创造了"真空市场"。}中国 AI 算力市场规模约 500 亿美元（英伟达 CEO 黄仁勋 2025 年语），且以每年 30\%+ 增速扩张。当进口渠道关闭，这一巨量需求转向国产方案，为本土企业提供了此前不可能获得的规模化部署机会。

\textbf{第二，政府投资形成"超级催化剂"。}国家集成电路产业投资基金（大基金）三期规模 3440 亿元（约 \$47B），2024 年成立；2026 年两会进一步释放约 \$70B 的 AI 与半导体补贴信号。政策资金不仅直接投向芯片公司，更催熟了代工（中芯国际 7nm N+2 制程）、封装、设备等全产业链。

\textbf{第三，客户"被迫试用"打破了生态锁定。}过去国内企业即使芯片性能达标，也难以撬动已深度绑定 CUDA 生态的客户。出口管制让客户别无选择，被迫适配国产方案，反而加速了软件生态的成熟——寒武纪 Cambricon Neuware、摩尔线程 MUSA（CUDA 兼容层）、华为 CANN 的用户基数均在 2025 年实现数量级增长。
\end{whybox}

%% ----------------------------------------------------------------
%%  Major Company Profiles
%% ----------------------------------------------------------------

\subsubsection{寒武纪：从"中国小 NVIDIA"到首次盈利}

寒武纪（Cambricon, SSE:688256）是中国 AI 芯片领域最具标志性的企业。创始人陈天石 1985 年出生，16 岁考入中科大少年班，博士毕业于中科院计算所，2016 年与中科算源共同创立寒武纪，寓意"智能大爆发"。公司于 2020 年 7 月登陆科创板，成为"AI 芯片第一股"，上市首日市值即突破千亿元。

然而，高光之下是漫长的亏损期。2017 至 2024 年，寒武纪连续 8 年亏损，累计超 54 亿元——巨额研发投入是主因，截至 2025 年 6 月，公司累计申请专利 2774 项，其中发明专利占比超 97\%。资本市场的耐心来自一个核心信念：国产 AI 芯片赛道终将兑现。

兑现时刻在 2025 年到来。2025 全年，寒武纪实现营业收入 64.97 亿元，同比增长 453.21\%；归母净利润 20.59 亿元，首次实现年度盈利，正式"摘 U"——自 2026 年 3 月 16 日起，股票简称从"寒武纪-U"变更为"寒武纪"。云端产品线（思元 590/690 系列芯片及加速卡）贡献了超过 99\% 的营收，在运营商、金融、互联网等行业实现规模化部署。

寒武纪的核心竞争力在于"芯片 + 基础软件"的全栈能力。其 Cambricon Neuware 平台兼容主流开源大模型，思元 590 在推荐系统场景中的能效比甚至达到 NVIDIA H100 的 1.8 倍。2025 年，公司以 1195.02 元/股完成定增，募资 39.85 亿元用于大模型芯片研发，下一代思元 790 计划于 2026 年推出。

2025 年《胡润全球富豪榜》上，陈天石以 870 亿元身家登顶中国 AI 行业首富；一年后财富飙至 1750 亿元，胡润董事长评价其为"未来中国首富的有力候选人"。寒武纪市值曾一度达到 4630 亿元，在科创板 AI 芯片板块中居首。东吴证券将其 2025--2027 年营收预测上调至 67.71/135.35/230.04 亿元，维持"买入"评级。

\subsubsection{海光信息：唯一国产 x86 的双线战略}

海光信息（Hygon, SSE:688041）成立于 2014 年，其独特之处在于早年获得了 AMD 的 x86 架构授权——这使其成为全球除英特尔、AMD 之外极少数能合法生产 x86 服务器 CPU 的企业。公司采取"C86 CPU + DCU GPU"双产品线战略：C86 系列直接进入国产信创服务器市场，DCU（Deep Computing Unit）则定位 AI 训练与推理加速，对标 NVIDIA 的 GPU 产品线。

2025 年，海光信息实现营收 143.76 亿元，同比增长 57\%，净利润 25.4 亿元，是中国 AI 芯片板块中营收规模最大的企业。其核心优势在于"双轮驱动"——CPU 端受益于信创工程的确定性需求（政府、金融、电信等关键行业的国产替代），GPU 端则乘 AI 算力爆发之东风。中曙光（Sugon）是其最大客户与渠道合作伙伴，海光 DCU 已部署于全国多个智算中心。科创板市值排名中，海光信息以约 4711 亿元位列第二，仅次于寒武纪。

\subsubsection{地平线：车载 AI 芯片的"中国 Mobileye"}

地平线（Horizon Robotics, SEHK:9660）由前百度深度学习实验室（IDL）负责人余凯于 2015 年创立，2024 年 10 月在港交所上市，融资 54 亿港元。与前述企业聚焦数据中心不同，地平线选择了一条差异化赛道——车载 AI 芯片，以"征程"（Journey）系列处理器为核心产品。

2025 年，地平线营收达 37.6 亿元，同比增长 57.7\%；毛利润 24.3 亿元，毛利率维持在 64.5\% 高位，其中汽车业务毛利率高达 67.2\%。征程系列芯片全年出货超 400 万套，同比增长 39\%；更关键的是，支持城区/高速 NOA（Navigate on Autopilot）的中高阶芯片出货量达 180 万套，较上年增长近 5 倍，贡献了超 80\% 的产品与解决方案业务收入。

地平线的商业模式兼具"卖芯片"和"卖方案"双重属性。产品解决方案业务营收跃升至 16.2 亿元（+144.2\%），占比从 28\% 提升至 43\%；授权及服务业务 19.4 亿元（+17.4\%）。其 Horizon SuperDrive（HSD）方案已获 10 家 OEM 品牌累计 20 余款车型定点，首搭车型星途 ET5 上市 8 周即交付 2.5 万套，顶配版在该车型总销量中占比 83\%。大众集团旗下 CARIAD 与地平线成立合资公司，进一步验证了其技术在全球一线车企中的认可度。

2025 年中国乘用车市场智能辅助驾驶渗透率达 67.6\%，其中 NOA 功能渗透率从 21.6\% 飙升至 42.6\%。地平线以 47.7\% 的国产 ADAS 芯片市占率，牢牢占据"智驾平权"浪潮中的核心位置。

%% ----------------------------------------------------------------
%%  GPU Four Dragons + Other Companies
%% ----------------------------------------------------------------

\begin{keybox}[GPU 四小龙 IPO 数据汇总]
2025 年底至 2026 年初，被称为上海"GPU 四小龙"的四家企业密集登陆资本市场，标志着中国国产 GPU 行业从"PPT 造芯"阶段迈入量产落地与资本化的成熟期。

\begin{itemize}
  \item \textbf{摩尔线程（Moore Threads, SSE:688795）}——2020 年由前 NVIDIA 中国区总经理张建中创立，MUSA 架构主打 CUDA 兼容。2025 年 12 月科创板上市，融资 80 亿元，\textbf{首日涨幅 +425\%}。2025 年营收 15.06 亿元（+243\%）。
  \item \textbf{沐曦（MetaX, SSE:688802）}——2020 年由 ex-AMD 高级总监陈维良创立。2025 年 12 月科创板上市，融资 42 亿元，\textbf{首日涨幅 +693\%}。GPU 四小龙中技术路线最纯粹的一家。
  \item \textbf{壁仞科技（Biren, SEHK:6082）}——2019 年由前商汤科技总裁张文创立，BR100/BR20X GPGPU 系列。2026 年 1 月 2 日港交所上市，融资 55.8 亿港元，\textbf{首日涨幅 +76\%}。2023 年被列入 Entity List。
  \item \textbf{天数智芯（Iluvatar CoreX, SEHK:9903）}——2015 年成立，2018 年启动通用 GPU 设计。2026 年 1 月 8 日港交所上市，发行价 144.6 港元，IPO 市值约 354 亿港元。中国首家实现"训推双量产"的 GPGPU 企业。2024 年收入 5.4 亿元。
\end{itemize}

四家公司两家选择科创板、两家赴港，合计募资超 230 亿元。上海作为集聚超 1200 家集成电路企业、汇聚全国约 40\% 产业人才和近 50\% 创新资源的城市，为"四小龙"提供了独特的产业集群土壤。
\end{keybox}

\subsubsection{昆仑芯与燧原科技：大厂分拆与战略投资}

昆仑芯（Kunlunxin）2021 年从百度分拆独立，继承了百度自 2010 年代中期即开始的 AI 芯片自研积累。昆仑 I 代于 2019 年量产，是中国首款云端通用 AI 芯片；昆仑 II 代性能提升 2--3 倍，部署于百度搜索、自动驾驶等核心场景；昆仑 III 代于 2024 年发布。2024 年营收超 35 亿元。2026 年 1 月，昆仑芯向港交所递交 IPO 申请，目标融资 \$2B，一旦上市将进一步扩充国产 AI 芯片的资本军团。

燧原科技（Enflame）由 AMD 系资深芯片专家赵立东和张亚林于 2018 年在上海创立。赵立东 2007--2014 年在 AMD 任计算事业部高级总监，主导参与了 AMD 中国研发中心的成立。公司自研 GCU（General Compute Unit）架构，推出邃思（CloudBlazer）系列训练与推理芯片。2026 年 1 月 22 日，燧原科技科创板 IPO 获受理，拟募资 60 亿元。然而，作为"四小龙的老大哥"，燧原科技的身份存在争议——有行业媒体指出其产品属于 ASIC 专用芯片类别而非通用 GPU，因此严格意义上"GPU 四小龙"应为摩尔线程、沐曦、壁仞和天数智芯。

\begin{warnbox}[风险警示：客户集中度与 Entity List]
\textbf{客户集中度风险。}燧原科技 2025 年 71.84\% 的营收来自腾讯（持股 19.94\%）。当一家芯片公司的收入超过七成依赖单一客户时，其定价权、产品独立性和抗风险能力均令人担忧。类似地，寒武纪年报显示前五大客户频繁轮换——从终端 IP 授权到智能计算集群再到云端产品线，核心客户的不稳定性折射出国产芯片市场的早期特征。

\textbf{Entity List 影响。}壁仞科技于 2023 年被美国列入实体清单，限制其采购含美技术成分超 25\% 的产品，对先进制程代工和 EDA 工具获取构成直接制约。算能（Sophgo, 2019 年从比特大陆分拆）则于 2025 年 1 月被加入 Entity List，原因是 TSMC 代工的芯片疑似流入华为。被列入清单并不意味着公司"死亡"——壁仞 2026 年 1 月仍成功港股上市并首日暴涨 76\%，但供应链风险和技术迭代受限是长期隐忧。
\end{warnbox}

\subsubsection{算能：RISC-V 的异类路径}

算能（Sophgo）2019 年从比特大陆（Bitmain）分拆，创始人詹克团。与前述企业基于自研或 x86/ARM 架构不同，算能押注"TPU + RISC-V"双线战略。其 BM1684X 系列 TPU 面向 AI 推理，同时积极布局 RISC-V 开源指令集架构的 AI 处理器。2025 年，算能推出全球首台基于 RISC-V 架构的服务器并成功运行 DeepSeek 模型，验证了 RISC-V 在 AI 推理场景中的可行性。尽管 2025 年 1 月被列入 Entity List，算能在 RISC-V 赛道的布局使其在架构层面具备了一定程度的"去美化"优势——RISC-V 作为开源指令集，不受单一国家的许可控制。

%% ----------------------------------------------------------------
%%  Summary Table
%% ----------------------------------------------------------------

\begin{table}[htbp]
\centering
\caption{中国主要 AI 芯片公司全景（截至 2026 年 3 月）}
\label{tab:china-ai-chips}
\small
\begin{tabularx}{\textwidth}{@{}l l X r l l@{}}
\toprule
\rowcolor{figLightGray}
\textbf{公司} & \textbf{核心产品线} & \textbf{技术路线} & \textbf{2025 营收} & \textbf{上市状态} & \textbf{核心客户} \\
\midrule
寒武纪 & 思元 590/690 & 云端 ASIC & 64.97 亿 & SSE:688256 & 运营商/金融/互联网 \\
海光信息 & C86 CPU + DCU & x86 + GPU 双线 & 143.76 亿 & SSE:688041 & 中曙光/信创体系 \\
昆仑芯 & 昆仑 I/II/III & 百度自研架构 & $>$35 亿 & 港股 IPO 申请中 & 百度生态 \\
燧原科技 & 邃思系列 GCU & AI 专用芯片 & --- & 科创板 IPO 受理 & 腾讯（71.8\%） \\
壁仞科技 & BR100/BR20X & GPGPU & --- & SEHK:6082 & 政企/数据中心 \\
摩尔线程 & MUSA 架构 GPU & CUDA 兼容 GPU & 15.06 亿 & SSE:688795 & 互联网/云厂商 \\
沐曦 & 曦云/曦思系列 & 高性能 GPU & --- & SSE:688802 & 数据中心 \\
天数智芯 & 天垓 + 智铠 & 训推双线 GPGPU & 5.40 亿 & SEHK:9903 & 多行业 \\
地平线 & 征程系列 & 车载 AI SoC & 37.6 亿 & SEHK:9660 & 大众/比亚迪/理想 \\
算能 & BM1684X TPU & TPU + RISC-V & --- & 未上市 & 边缘推理/IoT \\
\bottomrule
\end{tabularx}
\end{table}

%% ----------------------------------------------------------------
%%  Structural Analysis
%% ----------------------------------------------------------------

\subsubsection{生态格局与竞争态势}

上述十家企业可沿两个维度划分。\textbf{纵轴}是应用场景：数据中心训练（寒武纪、海光、昆仑芯、燧原、四小龙）、数据中心推理（天数智芯、算能）、车载/边缘（地平线）。\textbf{横轴}是架构路线：ASIC 专用芯片（寒武纪、燧原）、通用 GPU/GPGPU（海光 DCU、壁仞、摩尔线程、沐曦、天数智芯）、x86 CPU（海光 C86）、车载 SoC（地平线征程）、RISC-V（算能）。

值得注意的是，这一生态正从"单点突破"向"体系竞争"演进。华为昇腾作为"一超"，正通过开放灵衢（LingQiu）互联协议和 CANN 软件栈构建开放生态，试图成为中国版的"NVIDIA CUDA + NVLink"。而"多强"阵营的企业则面临两难：一方面需要差异化以避免与华为正面竞争（如地平线专注车载、摩尔线程主打 CUDA 兼容），另一方面又必须融入华为主导的算力基础设施标准。

2026 年 1 月，平头哥（阿里巴巴半导体子公司）传出独立 IPO 消息，其 PPU 芯片性能接近 NVIDIA H20，背靠阿里 3800 亿元 AI 基建投资。若平头哥成功上市，"一超多强"格局可能演变为"双超多强"，进一步加速市场洗牌。

\subsubsection{从 IPO 狂欢到价值检验}

四小龙的 IPO 首日涨幅令人瞠目——摩尔线程 +425\%、沐曦 +693\%、壁仞 +76\%、天数智芯开盘 +31.5\%。但狂欢背后，多数企业仍处于大额亏损状态。天数智芯 2022--2024 年累计亏损超 22 亿元，研发费用常年超过营收。这些企业的估值逻辑建立在"国产替代的确定性"之上——投资者买入的不是当期利润，而是中国 AI 算力 500 亿美元市场中英伟达份额归零后的填补机会。

然而，当所有企业都瞄准同一块市场时，过度竞争与估值泡沫的风险同样不容忽视。科创板市值前四已被 AI 芯片公司霸榜——寒武纪 5427 亿元、海光信息 4711 亿元、摩尔线程 3317 亿元、沐曦 3149 亿元（2025 年 12 月 18 日数据）。这四家公司的合计市值超过 1.66 万亿元，而 2025 年合计营收约 224 亿元，隐含的 PS（市销率）超过 70 倍。即便考虑到高成长性，这一估值水平也意味着市场已将未来数年的国产替代红利充分定价，留给后来者的安全边际并不宽裕。

最终，决定这场"突围之路"成败的，不仅是芯片的晶体管数量或 TOPS 算力指标，更是软件生态的成熟度、客户的真实复购意愿、以及在中芯国际 7nm 良率约 30\% 条件下的量产交付能力。正如天数智芯在上市时所言："集体推进资本进程对整个行业绝对是利好，这标志着国产 GPU 行业已经跨越了早期的'PPT 造芯'阶段，真正迈向了量产落地和资本化的成熟阶段。"从 PPT 到量产，从亏损到盈利，从被封锁到自主可控——中国 AI 芯片的突围之路，才刚刚走过最艰难的第一程。


% §2.3 Training Infrastructure
% ch02-s3-training.tex  — §2.3 Training Infrastructure
% Part of Chapter 2: Infrastructure Layer
% Input from: ch02-infrastructure.tex (or main file)

\section{训练基础设施}
\label{sec:infra-training}

Berkeley 黑帮：一个实验室如何批量制造独角兽。

加州大学伯克利分校的 RISELab（及其前身 AMPLab）
也许是人类历史上商业转化率最高的计算机科学实验室。
从这里走出了 Apache Spark $\to$ Databricks（估值 \$620 亿），
Ray $\to$ Anyscale（估值 \$10 亿），
以及最新的 vLLM $\to$ Inferact（种子轮 \$1.5 亿，估值 \$8 亿）。
同一套剧本反复上演：
\textbf{先开源一个足够好的系统级项目，等它成为事实标准，
然后以商业公司的形式提供企业版和托管服务}。
这不是偶然——Berkeley 的文化把``Impact''定义为
``被百万开发者使用''而非``被 ICML 接收''，
而硅谷的风投生态恰好能将这种 Impact 货币化。

训练基础设施赛道是这一模式的集中展演。
从分布式计算框架到实验跟踪平台，
从 GPU 编排到高效训练算法，
这里的每一家公司都在试图回答同一个问题：
\textbf{如何让训练一个大模型的过程变得更快、更便宜、更可重复？}
而答案的结局，往往不是独立 IPO——而是被巨头收购。

\begin{corebox}[Berkeley-to-Startup 流水线]
Berkeley 计算机科学系在 AI 基础设施领域形成了一条独特的
``开源→标准化→商业化''流水线：

\begin{enumerate}[noitemsep]
  \item \textbf{学术孵化}——在 RISELab / Sky Computing Lab 中
        开发系统级开源项目，发表顶会论文建立学术声誉
  \item \textbf{社区扩散}——以 Apache / Linux Foundation 等中立基金会
        托管项目，吸引产业界贡献者，形成事实标准
  \item \textbf{商业剥离}——创始团队成立公司，提供企业版、
        托管云服务、SLA 支持等增值层
  \item \textbf{生态锁定}——通过 API 兼容性和社区惯性
        建立竞争壁垒，最终实现 IPO 或被战略收购
\end{enumerate}

\medskip
\small
\begin{tabularx}{\textwidth}{@{}l l l r l@{}}
\toprule
\rowcolor{figLightGray}
\textbf{开源项目} & \textbf{商业公司} & \textbf{创始人} &
  \textbf{估值} & \textbf{状态} \\
\midrule
Apache Spark  & Databricks   & Ion Stoica 等  & \$620B & 独立（待 IPO）\\
Ray           & Anyscale     & Ion Stoica 等  & \$1B   & 独立 \\
vLLM          & Inferact     & Simon Mo 等    & \$800M & 种子轮 \\
Apache Kafka  & Confluent    & Jay Kreps 等   & 上市   & NASDAQ: CFLT \\
\bottomrule
\end{tabularx}
\end{corebox}

然而，训练基础设施赛道的另一面是\textbf{大规模整合}。
在本节覆盖的八家公司中，五家已被收购：
MosaicML $\to$ Databricks，
Determined AI $\to$ HPE，
Run:ai $\to$ NVIDIA，
Weights \& Biases $\to$ CoreWeave，
加上 MLOps 赛道外围的无数沉默者。
这一整合浪潮的逻辑清晰：
\textbf{当 AI 基础设施的买家只有少数几家超大规模平台公司时，
``成为生态的一部分''比``独立上市''更现实}。

\begin{keybox}[训练基础设施赛道收购汇总]
\small
\begin{tabularx}{\textwidth}{@{}l l r l l@{}}
\toprule
\rowcolor{figLightGray}
\textbf{被收购方} & \textbf{收购方} & \textbf{价格} &
  \textbf{时间} & \textbf{战略逻辑} \\
\midrule
MosaicML       & Databricks  & \$1.3B  & 2023/6  & 补全模型训练能力 \\
Determined AI  & HPE         & 未披露   & 2021/6  & 集成 Cray HPC 生态 \\
Run:ai         & NVIDIA      & \$700M  & 2024/12 & GPU 编排 + 软件栈 \\
Weights \& Biases & CoreWeave & \$1.7B & 2025/5 & 构建端到端 AI 云 \\
\bottomrule
\end{tabularx}
\end{keybox}

以下逐一分析八家核心公司的技术路线、商业模式与命运走向。


\subsection{Anyscale：Ray 背后的商业引擎}

Ion Stoica 是那种让硅谷风投无法拒绝的连续创业者。
作为 Apache Spark 的联合创始人，
他在 2013 年将 Spark 的商业化公司 Databricks 带到了
如今 \$620 亿估值的高度。
2019 年，Stoica 再次出手，
围绕 Berkeley RISELab 孵化的分布式计算框架 Ray 创立了 Anyscale。

Ray 的设计哲学与 Spark 截然不同。
Spark 面向批处理数据分析（MapReduce 的继承者），
而 Ray 面向\textbf{通用分布式计算}——
特别是 AI 训练和推理中需要的异构、动态、有状态的工作负载。
一个 Ray 集群可以同时运行数据预处理（Ray Data）、
分布式训练（Ray Train）、超参搜索（Ray Tune）、
模型服务（Ray Serve）和强化学习（RLlib），
所有任务共享同一个资源调度器。

这种通用性让 Ray 成为了 AI 基础设施的``隐形冠军''。
截至 2026 年初，Ray 的 PyPI 累计下载量达 2.37 亿次，
GitHub 星标超过 3.6 万。
更关键的是其用户名单：
\textbf{OpenAI 使用 Ray 训练了 GPT-4 和后续模型}，
Anthropic、ByteDance、Uber、Spotify 等公司都将 Ray 作为
分布式 AI 工作负载的核心编排层。
Ray 还是 PyTorch Foundation 的成员项目，
与 PyTorch 的集成深度远超任何竞品。

商业化方面，Anyscale 提供托管的 Ray 平台（Anyscale Platform），
在 AWS 和 GCP 上提供一键部署、自动扩缩容和企业级监控。
公司累计融资约 \$2.6 亿，估值 \$10 亿。
相比 Databricks 的火箭式增长，
Anyscale 的商业化步伐显得更为克制——
这部分因为 Ray 的开源版本已经足够强大，
用户自行部署的门槛并不高。
如何在``开源足够好''和``企业版值得付费''之间找到平衡，
是 Anyscale 面临的核心挑战
（参见 \textit{Emergence} 第五章对分布式训练编排的技术讨论）。


\subsection{Lightning AI：从 PyTorch 封装到全栈 AI 云}

William Falcon 在 2019 年创立 Lightning AI 时，
最初只是想解决一个让每个深度学习研究者头疼的问题：
\textbf{PyTorch 代码的工程化太痛苦了}。
PyTorch Lightning 的核心设计是将训练循环中的
``科学逻辑''（模型、损失函数、优化器）
与``工程逻辑''（分布式、混合精度、检查点、日志）彻底分离，
让研究者只需写前者，框架自动处理后者。

这个看似简单的抽象层一炮而红：
PyTorch Lightning 在 GitHub 上获得超过 2.6 万颗星，
PyPI 累计下载量超过 1.6 亿次，
成为仅次于 PyTorch 本身的最流行深度学习训练框架。

但真正的转折发生在 2026 年 1 月。
Lightning AI 宣布与 GPU 云服务商 Voltage Park 完成合并，
合并后估值达 \$25 亿，ARR 从 2024 年的 \$1800 万
飙升至超过 \$5 亿——
\textbf{不到两年增长近 28 倍}。
合并后的公司拥有超过 40 万名开发者用户，
提供从 GPU 集群到训练框架到模型部署的全栈 AI 云服务。

这次合并揭示了训练基础设施赛道的一个深层趋势：
\textbf{纯软件的训练工具正在向``软件 + 算力''一体化演进}。
当用户在 Lightning 平台上编写训练代码时，
平台可以自动分配 Voltage Park 的 GPU 资源，
实现从代码到集群的无缝衔接。
这种垂直整合策略让 Lightning AI 从一个开源框架公司，
蜕变为 AWS / GCP 在 AI 训练领域的垂直竞争者。


\subsection{MosaicML：被 Databricks \$13 亿收购的效率先锋}

Naveen Rao 的创业履历本身就是 AI 基础设施赛道的缩影。
2014 年创立神经形态芯片公司 Nervana Systems，
2016 年被 Intel 以 \$3.34 亿收购；
在 Intel 担任 AI 负责人后，
2020 年再次创业成立 MosaicML。

MosaicML 的核心竞争力是\textbf{训练效率}。
其开源库 Composer 集成了数十种训练加速技术
（MixUp、Label Smoothing、Sharpness-Aware Minimization 等），
让用户可以像搭积木一样组合不同的``训练配方''，
在不改变模型架构的前提下将训练速度提升数倍。
更引人注目的是 MosaicML 自研的 MPT（MosaicML Pretrained Transformer）系列模型——
这些模型证明了一个关键论点：
\textbf{通过优化训练流程，中等规模的团队也能训练出
与大厂模型质量相当的大语言模型}。

2023 年 6 月，Databricks 以 \$13 亿收购 MosaicML。
这笔交易的战略逻辑清晰：
Databricks 拥有全球最大的企业数据平台，
但缺乏模型训练能力；
MosaicML 拥有最先进的训练技术，
但缺乏数据和分发渠道。
收购后，MosaicML 的技术被整合为 Databricks 的 Foundation Model Training 服务，
让企业客户可以在自己的数据湖上直接训练定制化大模型。

故事的后续同样精彩：
Rao 在 2025 年 9 月离开 Databricks，
创办了 Unconventional AI——
一家瞄准类脑计算（neuromorphic computing）的新公司，
\textbf{种子轮即融资 \$4.75 亿，估值 \$45 亿}。
这个天文数字的种子轮
不仅刷新了 AI 领域的融资纪录，
更说明了一个残酷的事实：
在 AI 基础设施赛道，
\textbf{创始人的连续创业信用比任何技术路线都值钱}。


\subsection{Determined AI：被 HPE 收入 HPC 版图}

Determined AI 的故事相对低调但颇具代表性。
公司由 Neil Conway、Evan Sparks 和 CMU 教授 Ameet Talwalkar
于 2017 年在旧金山创立，
开发了一套开源的分布式训练平台，
核心功能包括自动分布式训练、超参搜索、实验管理和资源调度。

2021 年 6 月，HPE 收购 Determined AI
（价格未披露，估计在数千万美元级别），
将其技术集成到 HPE 的 Cray 超算生态中。
对于 HPE 而言，这是一笔典型的``补短板''收购：
HPE 拥有世界顶级的 HPC 硬件（Cray 超算、Slingshot 互联），
但缺乏面向 AI 训练的现代软件栈。
Determined AI 的平台恰好填补了这一空白，
让 HPC 客户（国家实验室、大学、国防机构）
可以在 Cray 集群上运行 PyTorch 训练任务，
而无需从零搭建分布式训练基础设施。

Determined AI 的案例说明了一条被低估的退出路径：
\textbf{不是每家 AI 基础设施创业公司都需要追逐十亿美元级别的结局。
在正确的时间被正确的买家收购，
同样可以让技术找到最大化应用的归宿}。


\subsection{Run:ai：GPU 编排之王归于 NVIDIA}

当企业开始大规模部署 GPU 集群时，
一个隐藏的问题迅速浮出水面：
\textbf{GPU 的平均利用率低得惊人——通常不到 30\%}。
数据科学家提交训练任务后独占整组 GPU，
即使在数据加载或评估阶段 GPU 大部分时间处于空闲，
资源也无法被其他团队复用。

Run:ai（2018 年，Tel Aviv）正是瞄准了这一痛点。
其核心技术是\textbf{GPU 虚拟化与动态编排}：
将物理 GPU 抽象为可分割、可共享的虚拟资源，
并通过 Kubernetes 原生的调度器实现跨团队、跨项目的弹性分配。
一块 A100 可以同时服务于两个不同的训练任务，
或者在一个任务空闲时自动将算力释放给另一个任务。

Run:ai 累计融资约 \$1.18 亿，客户涵盖金融、医疗、科研等行业。
2024 年 4 月，NVIDIA 宣布收购 Run:ai；
经过欧盟反垄断审查后，交易于 2024 年 12 月以约 \$7 亿完成。
收购完成后，NVIDIA 做出了一个出人意料的决定：
\textbf{将 Run:ai 的核心软件开源}。

这一策略的逻辑耐人寻味。
开源 Run:ai 意味着任何 GPU 集群（包括竞品芯片）
都可以使用这套编排软件，
但实际效果是进一步巩固了 NVIDIA 作为
``AI 基础设施事实标准''的地位——
正如 Android 之于 Google，
免费的软件层反而加深了硬件生态的锁定。


\subsection{Weights \& Biases：MLOps 标杆被 CoreWeave 收编}

在所有训练基础设施公司中，
Weights \& Biases（W\&B）也许是\textbf{开发者口碑最好的一家}。
创始人 Lukas Biewald 在 2017 年将这家公司定位为
``AI 团队的 GitHub''——
一个用于实验跟踪、模型版本管理、数据可视化和协作的平台。

W\&B 的成功在于产品的\textbf{极致简洁}。
只需在训练脚本中加入两行代码（\texttt{import wandb; wandb.init()}），
就可以自动记录所有超参数、指标、系统状态和输出文件。
这种近乎零摩擦的集成体验让 W\&B 在 AI 研究社区迅速传播，
积累了超过 10 万名客户，
包括 OpenAI、NVIDIA、Microsoft、Toyota 等。
公司累计融资约 \$2.5 亿，估值达 \$12.5 亿。

2025 年 3 月，CoreWeave 宣布以约 \$17 亿收购 W\&B。
CoreWeave 是 NVIDIA GPU 云的最大独立运营商，
但其产品栈一直停留在``卖算力''层面。
收购 W\&B 让 CoreWeave 一举获得了
从实验管理到模型训练到部署监控的\textbf{完整软件层}，
实现了从``GPU 云''到``AI 开发平台''的跃迁。

对于 MLOps 赛道而言，这笔交易传递了一个清晰信号：
\textbf{独立 MLOps 工具的天花板已经可见}。
当算力提供商开始纵向整合软件栈时，
纯工具层公司面临的选择只有两个：
向上扩展成平台（如 Lightning AI + Voltage Park），
或者被平台收购（如 W\&B $\to$ CoreWeave）。


\subsection{ClearML：精益生存的开源 MLOps}

并非每家 AI 基础设施公司都需要数亿美元的融资。
ClearML（2016 年，Israel）证明了一条截然不同的路径：
\textbf{用极低的资本消耗，服务足够多的用户}。

ClearML 的产品是一套开源的端到端 MLOps 平台，
覆盖实验管理、数据版本化、流水线编排和模型部署。
与 W\&B 的 SaaS 优先策略不同，
ClearML 坚持\textbf{开源优先 + 自托管为主}的路线：
企业客户可以在自己的服务器上完整部署 ClearML，
无需将任何数据发送到外部。
这一策略在对数据主权敏感的行业（金融、国防、医疗）中尤具吸引力。

公司累计融资仅约 \$1600 万——
不到 W\&B 的十五分之一——
但积累了超过 25 万名用户。
截至 2026 年 3 月，ClearML 仍然保持独立运营，
并在 GTC 2026 上发布了 Platform Management Center，
为企业 IT 管理员提供 AI 基础设施的财务透明度和用量分析。

ClearML 的生存策略提出了一个值得深思的问题：
在 VC 驱动的 AI 创业生态中，
\textbf{``精益''是否也是一种竞争优势？}
当那些融资数亿的竞争对手纷纷被收购或关停时，
ClearML 以极低的烧钱速率持续迭代，
反而获得了独立生存的自由度。


\subsection{Comet ML：数据驱动的 ML 生命周期管理}

Comet ML（2017 年，纽约）定位为 ML 全生命周期管理平台，
产品覆盖从实验跟踪到生产模型监控的完整链路。
与 W\&B 正面竞争，
Comet 的差异化在于更强的\textbf{生产环境监控}能力——
不仅跟踪训练阶段的实验指标，
还持续监控部署后模型的数据漂移、性能退化和公平性指标。

公司累计融资约 \$6300 万，团队仅 88 人。
2024 年营收达 \$1700 万，同比增长 66\%——
在 MLOps 赛道普遍增长放缓的背景下，这是一个不俗的成绩。
但 Comet 面临的挑战同样严峻：
W\&B 被 CoreWeave 收购后获得了算力平台的加持，
Databricks 通过收购 MosaicML 内建了 MLflow + 训练的一体化体验，
MLflow 本身作为开源项目又在不断蚕食商业 MLOps 工具的市场空间。

\textbf{夹在开源免费和巨头捆绑之间的独立 MLOps 公司，
生存空间正在被双向挤压。}
Comet 的 \$1700 万营收虽然在增长，
但与 Lightning AI 合并后的 \$5 亿 ARR 相比，
量级差距已经难以弥合。


\begin{whybox}[为什么训练基础设施赛道被大规模整合？]
在八家公司中，五家已被收购。这并非巧合，而是结构性必然：

\begin{enumerate}[noitemsep]
  \item \textbf{买家集中度极高}——
        AI 训练基础设施的终端客户集中在少数超大规模平台
        （云厂商、GPU 云、大模型公司），
        这些平台有强烈动机通过纵向整合降低自身客户的迁移成本
  \item \textbf{开源侵蚀商业护城河}——
        训练基础设施的核心技术大多以开源形式存在
        （Ray、PyTorch Lightning、MLflow、Determined AI），
        商业化增值层的差异化空间有限
  \item \textbf{``算力 + 软件''一体化趋势}——
        客户不想分别采购 GPU 集群和训练工具，
        而是希望获得一站式体验
        （CoreWeave + W\&B，Lightning + Voltage Park）
  \item \textbf{人才价值超过产品价值}——
        收购 MosaicML 的 \$13 亿中，
        很大一部分是为 Naveen Rao 团队的训练专家支付的``人才溢价''
\end{enumerate}

对于仍然独立的公司（Anyscale、ClearML、Comet），
问题不是``是否会被整合''，而是``何时、被谁、以什么价格''。
\end{whybox}

\subsubsection{训练基础设施赛道全景对比}

表~\ref{tab:training-infra-comparison} 汇总了本节覆盖的八家公司的关键指标。

\begin{table}[htbp]
\centering
\caption{训练基础设施创业公司对比（截至 2026 年 3 月）}
\label{tab:training-infra-comparison}
\small
\begin{tabularx}{\textwidth}{@{}l l l r l l@{}}
\toprule
\rowcolor{figLightGray}
\textbf{公司} & \textbf{成立} & \textbf{核心产品} &
  \textbf{总融资} & \textbf{状态} & \textbf{收购方 / 价格} \\
\midrule
Anyscale        & 2019 & Ray 分布式计算        & \$260M  & 独立      & --- \\
Lightning AI    & 2019 & PyTorch Lightning     & 合并    & 合并      & + Voltage Park / \$2.5B 估值 \\
MosaicML        & 2020 & Composer + MPT        & ---     & 被收购    & Databricks / \$1.3B \\
Determined AI   & 2017 & 分布式训练平台         & ---     & 被收购    & HPE / 未披露 \\
Run:ai          & 2018 & GPU 编排虚拟化        & \$118M  & 被收购    & NVIDIA / \$700M \\
W\&B            & 2017 & 实验跟踪 + MLOps      & \$250M  & 被收购    & CoreWeave / \$1.7B \\
ClearML         & 2016 & 开源 MLOps 平台       & \$16M   & 独立      & --- \\
Comet ML        & 2017 & ML 生命周期管理       & \$63M   & 独立      & --- \\
\bottomrule
\end{tabularx}
\end{table}

训练基础设施赛道的终局图景正在变得清晰：
\textbf{独立的纯软件工具层正在消融，
取而代之的是``算力 + 框架 + 工具''的垂直整合平台}。
CoreWeave 收购 W\&B、Lightning AI 与 Voltage Park 合并、
NVIDIA 开源 Run:ai——
每一个动作都指向同一个方向：
AI 训练的用户体验将越来越像
``打开浏览器，写代码，点击训练''，
而底层的分布式调度、GPU 编排、实验管理将全部隐入平台。

对于仍然独立的 Anyscale、ClearML 和 Comet 而言，
窗口期不会太长。
Berkeley 流水线的下一个产品 Inferact（vLLM 商业化）
已经以 \$8 亿的种子轮估值入场，
新一轮的开源→标准化→商业化循环即将开始。
在这个赛道上，唯一不变的是变化本身——
\textbf{而推动变化的引擎，
仍然是 Berkeley Soda Hall 那几间不起眼的实验室}。

% ===========================================================
\subsection{更多训练工具与框架创业公司}
\label{subsec:training-more}
% ===========================================================

除了上述八家核心公司，训练基础设施赛道还有一批
聚焦特定技术层面的小型创业公司和开源项目，
它们虽未达到独角兽规模，但在各自的技术缝隙中发挥着重要作用。

\subsubsection{HPC-AI Tech（ColossalAI）}
HPC-AI Tech（2021 年，Singapore）由新加坡国立大学教授
杨优（Yang You）创立，旗下开源项目 ColossalAI
（GitHub 41,000+ 星，190+ 贡献者）
是\textbf{大模型高效训练}领域最活跃的开源框架之一。
ColossalAI 提供数据并行、模型并行、流水线并行和异构训练的
一站式解决方案，核心优势是让中小团队也能在有限 GPU 上
训练和微调大模型。2024 年完成 \$5000 万 Series A 融资
（累计 \$5600 万），同时推出 Open-Sora 开源视频生成平台。
ColossalAI 对标的是 NVIDIA 的 Megatron-LM 和 Microsoft 的
DeepSpeed——但作为独立创业公司而非大厂内部项目，
它在商业化路径上面临更大挑战。

\subsubsection{Predibase（Ludwig → LoRAX → Rubrik 收购）}
Predibase（San Francisco）由 Apple 和 Uber 内部 AI 平台的
核心工程师创立，最初基于开源框架 Ludwig（声明式 ML 训练工具）
构建微调平台，后推出 LoRAX——
\textbf{首个支持同时服务数千个 LoRA 适配器的推理引擎}。
2025 年 6 月完成 \$2800 万 Series A 融资后，
\textbf{同月即被数据安全公司 Rubrik 以 \$1--5 亿收购}。
Predibase 的故事浓缩了一条典型路径：
开源项目 → 商业化 → 被垂直行业巨头收购以增强 AI 能力。
Rubrik 看中的是 Predibase 的微调技术——
让企业在自有数据上定制模型，同时保持数据安全。

\subsubsection{SkyPilot（Berkeley Sky Computing Lab）}
SkyPilot 是 Berkeley Sky Computing Lab 孵化的开源项目
（GitHub 7,000+ 星），由 UC Berkeley 教授 Ion Stoica
（Databricks/Anyscale 创始人）参与指导。
SkyPilot 的核心功能是\textbf{跨云 AI 工作负载编排}——
用一个统一接口在 AWS、GCP、Azure、Lambda、RunPod 等
任意云上启动训练或推理任务，自动选择最便宜或最快的可用 GPU。
对于多云环境下的 AI 团队，SkyPilot 解决了一个现实问题：
\textbf{GPU 在哪里便宜就去哪里跑}，无需手动管理多个云账号。
虽然尚未正式成立商业公司，
但 SkyPilot 是 Berkeley-to-Startup 流水线上的下一个潜在候选。


% §2.4 Inference Optimization and Serving
% =============================================================
% §2.4  Inference Optimization and Serving
% Book: The AI Startup Landscape
% =============================================================

\section{推理优化与服务}
\label{sec:infra-inference}

如果说训练是AI的研发实验室，那么推理就是AI的工厂车间。一个模型训练一次，却要被调用数十亿次。当ChatGPT在2023年创下用户增长纪录时，人们才意识到一个被长期忽视的事实：\textbf{运行模型的成本远超训练模型的成本}。据行业估算，到2025年，全球AI算力中约三分之二用于推理而非训练——而这个比例仍在快速攀升。

推理优化赛道的创业图景由此展开。从2022年到2026年，一系列由顶尖学术人才创立的公司——前Meta PyTorch团队、Caffe发明者、FlashAttention作者、LLVM/Swift之父——围绕一个共同目标展开竞争：\textbf{让每一个token的生成成本趋近于零}。这个赛道在短短三年内上演了万倍定价压缩、史上最大种子轮、NVIDIA系统性收购、以及一年三轮的融资奇观。

\begin{keybox}[推理API定价战：从奢侈品到大宗商品（2022--2026）]
\small
\begin{itemize}[noitemsep]
  \item \textbf{2022年末}：OpenAI GPT-3.5 Turbo 定价 \$2.00/M tokens（输入），推理是少数公司才能负担的奢侈品。
  \item \textbf{2023年中}：GPT-4 发布，定价 \$30.00/M tokens（输入）。但竞争者开始涌现，开源模型初露锋芒。
  \item \textbf{2024年初}：Mixtral 8x7B 掀起MoE浪潮，Together AI、Fireworks AI将GPT-3.5等效推理降至 \$0.20/M tokens。
  \item \textbf{2024年末}：DeepSeek-V2发布，中国市场引爆"百模大战"的定价内卷，\$0.14/M tokens。
  \item \textbf{2025年中}：GPT-4等效推理（以Llama 3.1 70B为代表）降至 \$0.20--0.40/M tokens，较2023年下降约100倍。
  \item \textbf{2026年初}：DeepSeek-R1的推理成本进一步被各平台压至 \$0.15/M tokens，小模型（8B）低至 \$0.03/M tokens。\textbf{三年间，同等能力的推理成本下降了约1000倍}。
\end{itemize}
\end{keybox}

这场定价战的驱动力来自两端：供给侧是推理引擎（vLLM、SGLang、TensorRT-LLM）的持续优化将硬件利用率提升了5--10倍；需求侧是AI Agent和Agentic Workflow的兴起，使单次用户交互从几十个token暴增至数万个token。推理不再是"调一次API"的事——一个AI Agent完成一项任务可能需要数十次模型调用、数百万token的推理量。推理优化的经济价值从"节省成本"升级为"使新范式成为可能"。

\begin{corebox}[推理优化技术栈：四大支柱]
现代LLM推理优化的核心技术形成了一个相互叠加的四层栈（详见\textit{Emergence}第九章的技术深度分析）：
\begin{enumerate}[noitemsep]
  \item \textbf{PagedAttention 与内存管理}。vLLM团队提出的PagedAttention将操作系统的虚拟内存分页思想引入KV Cache管理，解决了传统推理中因预分配连续内存导致的50--70\%显存浪费问题。这一技术让单GPU能同时服务的并发请求数提升了2--4倍，成为现代推理引擎的标配。
  \item \textbf{量化（Quantization）}。将模型权重从FP16/BF16压缩至FP8、INT4甚至更低精度。FP8量化在H100/H200上几乎无精度损失，可将推理速度提升近2倍。INT4量化（AWQ/GPTQ）在可接受的精度代价下实现约4倍加速，是成本敏感场景的首选。
  \item \textbf{Speculative Decoding}。用小模型（draft model）快速生成候选token序列，再用大模型一次性验证。在接受率较高的场景下（如代码补全），可以实现2--3倍的延迟改善，且\textbf{数学上保证输出分布与大模型完全一致}。
  \item \textbf{前缀缓存（Prefix Caching / RadixAttention）}。对共享系统提示或多轮对话中重复出现的前缀，缓存其KV值以避免重复计算。SGLang的RadixAttention将前缀组织为基数树结构，在多轮Agent场景下可节省30--50\%的计算开销。
\end{enumerate}
\end{corebox}

% ===========================================================
\subsection{Fireworks AI：推理云的企业级野心}
\label{subsec:fireworks}

Fireworks AI的创始人Lin Qiao是前Meta PyTorch基础设施团队的负责人——正是她的团队让PyTorch从一个研究框架成长为全球最流行的深度学习平台。2022年，她与同样来自Meta的联合创始人在Redwood City创立了Fireworks AI，赌注是：\textbf{训练框架的战争已经结束（PyTorch赢了），推理框架的战争才刚开始}。

这个判断被证明极具前瞻性。Fireworks AI从一开始就瞄准企业级推理市场，定位不是做"更便宜的API"，而是做\textbf{推理云}——一个让企业能在自己的数据安全边界内、以最优性能运行任意模型的平台。这种定位让Fireworks避开了与OpenAI的直接竞争，转而服务那些需要部署开源模型、定制微调、且对延迟和合规有严格要求的企业客户。

2025年10月，Fireworks AI完成了2.5亿美元的C轮融资，估值40亿美元，由Lightspeed和Index Ventures领投。公司披露年化经常性收入（ARR）约1.3亿美元，较前一年增长约20倍——这是推理赛道中罕见的、以实际收入支撑的高估值。客户名单包括Uber、Cursor、Notion、Shopify和Genspark，覆盖了从出行到开发工具到生产力的多个高价值垂直领域。

Fireworks的技术差异化在于其\textbf{复合AI推理引擎}（Compound AI Inference）。在实际的AI Agent应用中，一次用户请求往往需要编排多个模型——路由器模型判断意图、检索模型召回文档、推理模型生成回答、审核模型检查安全性。Fireworks的平台将这种多模型编排的调度、缓存和流量管理抽象为一层统一的基础设施，让企业不需要自己搭建复杂的推理管线。

Lin Qiao在C轮公告中写道："我们的客户告诉我们，他们不只是在买推理，他们在买\textbf{AI运行时}。"这个定位的精妙之处在于：API定价战只能压缩利润，但\textbf{运行时}是一个平台级别的价值主张——一旦企业将其AI基础设施建立在Fireworks上，迁移成本就会很高。

% ===========================================================
\subsection{Together AI：学术全明星的全栈雄心}
\label{subsec:together-inference}

Together AI的创始阵容是推理赛道中学术背景最华丽的：CEO Vipul Ved Prakash是连续创业者，联合创始人Tri Dao发明了FlashAttention（可能是2023--2025年间影响最广泛的单一推理优化技术），Christopher Ré和Percy Liang是斯坦福计算机系的知名教授。这支团队不仅理解推理优化的理论极限，更有将学术成果产品化的工程能力。

2022年成立于旧金山的Together AI，在2025年完成了3.05亿美元B轮融资，估值33亿美元。到2026年3月，公司正在洽谈以\textbf{75亿美元估值融资10亿美元}（The Information, 2026年3月5日），估值较B轮翻倍有余。平台已服务超过45万名开发者，年化营收据报突破10亿美元。

Together AI的定位是\textbf{全栈AI云}——覆盖推理、微调、数据处理的完整工作流。与Fireworks侧重企业推理不同，Together更强调对\textbf{开源生态}的深度投入：公司积极参与开源模型的训练和优化（如RedPajama数据集），Tri Dao的FlashAttention本身就是开源贡献，这种"开源优先"的策略为Together在开发者社区中建立了远超商业关系的信任纽带。

Together AI的技术护城河建立在三个层面：首先是\textbf{FlashAttention}的持续演进（从FlashAttention到FlashAttention-2到FlashAttention-3），确保在注意力计算效率上始终领先；其次是与NVIDIA的深度合作关系——NVIDIA是Together的重要投资方和云合作伙伴；最后是其\textbf{推理+微调+训练的闭环}，让企业可以在同一平台上完成从数据准备到模型部署的全流程，降低切换不同供应商的摩擦。

Together AI面临的核心挑战是：当推理成为一个低利润率的基础设施业务时，75亿美元的估值需要什么量级的营收来支撑？答案很可能在于\textbf{向上游延伸}——从卖推理到卖定制模型训练服务，从卖API到卖完整的AI基础设施栈。

% ===========================================================
\subsection{vLLM / Inferact：开源推理引擎的商业化之路}
\label{subsec:vllm}

vLLM的故事是开源基础设施如何从学术项目蜕变为数十亿美元商业机会的典型案例。

2023年，UC Berkeley的Woosuk Kwon等人发表了PagedAttention论文，并开源了基于这一技术的推理引擎vLLM。PagedAttention的核心洞察借鉴自操作系统：传统推理框架为每个请求预分配连续的KV Cache内存，导致大量内存碎片和浪费；PagedAttention将KV Cache分割为固定大小的"页"，按需分配和回收，就像操作系统管理虚拟内存一样。这一看似简单的工程抽象，将推理引擎的内存利用率从30--50\%提升至90\%以上。

vLLM迅速成为开源LLM推理的\textbf{事实标准}。截至2026年初，项目已累计超过\textbf{40,000颗GitHub星}，被部署在超过\textbf{40万块GPU}上，几乎所有主流推理服务——Fireworks AI、Together AI、Baseten、Anyscale——都在底层使用vLLM或其衍生版本。

2026年1月，vLLM的核心开发者以\textbf{Inferact}的名义正式成立商业公司，并宣布完成了\textbf{1.5亿美元种子轮融资}，估值8亿美元，由a16z领投、Lightspeed跟投——这是\textbf{史上最大规模的种子轮融资之一}（TechCrunch, 2026年1月22日）。这个数字的震撼之处在于：这不是一个有营收的成熟公司，而是一个刚刚从开源项目"spin out"的初创团队，纯粹凭借技术影响力和市场地位获得了这样的估值。

Inferact面临的核心悖论是"开源商业化困境"：vLLM之所以成功，正是因为它是免费和开源的；但一旦商业化，就必须在开源社区价值和商业收入之间找到平衡。Inferact的策略是效仿Red Hat模式——开源引擎保持免费，围绕其提供企业级的托管服务、技术支持、性能优化和安全审计。关键问题是：当你的竞争对手（Fireworks、Together）也在使用你的开源代码时，你的"企业版"能提供多少额外价值？

\begin{table}[htbp]
\centering
\caption{主流推理引擎性能对比（2026年初，Llama 3.1 70B，单节点8$\times$H200）}
\label{tab:inference-engines}
\small
\begin{tabularx}{\textwidth}{l r r r X}
\toprule
\rowcolor{figLightGray}
\textbf{引擎} & \textbf{吞吐量 (tok/s)} & \textbf{TTFT (ms)} & \textbf{部署难度} & \textbf{特色} \\
\midrule
vLLM 0.7+        & $\sim$5,800 & $\sim$85  & 低   & PagedAttention; 社区最大; 插件生态丰富 \\
SGLang 0.4+      & $\sim$6,200 & $\sim$78  & 中   & RadixAttention; 多轮Agent场景领先; 结构化输出 \\
TensorRT-LLM     & $\sim$7,000 & $\sim$70  & 高   & NVIDIA官方; 极致性能; 需手动编译优化 \\
TGI (HuggingFace) & $\sim$4,200 & $\sim$110 & 低   & 开箱即用; HuggingFace生态集成; 适合快速原型 \\
\bottomrule
\end{tabularx}

\medskip
\noindent\small 注：吞吐量为批处理场景下测量值，实际性能因batch size、序列长度、量化策略等因素差异较大。数据综合多个独立基准测试报告（Anyscale, 2025; ArtificialAnalysis, 2026）。
\end{table}

% ===========================================================
\subsection{SGLang / RadixArk：Agent推理的另一条路线}
\label{subsec:sglang}

如果说vLLM代表了"通用推理引擎"的路线，SGLang则选择了一条更垂直的道路——\textbf{为AI Agent和复杂推理工作流优化}。

SGLang同样诞生于UC Berkeley，核心开发者Ying Sheng曾在xAI工作（参与了Grok模型的推理基础设施建设）。SGLang的核心技术是\textbf{RadixAttention}——一种将多轮对话和Agent调用中的前缀KV Cache组织为基数树（Radix Tree）的数据结构。在典型的AI Agent场景中，系统提示和对话历史在多次模型调用中大量重复；RadixAttention自动识别和复用这些重复前缀，避免冗余计算，在多轮Agent场景下可节省30--50\%的推理开销。

2026年1月，SGLang项目以\textbf{RadixArk}的名义成立商业公司，由Accel领投，估值\textbf{4亿美元}（TechCrunch, 2026年1月21日）。虽然估值不及Inferact的8亿美元，但RadixArk的差异化定位——聚焦Agent推理——可能在AI Agent爆发的2026年找到更精准的PMF（Product-Market Fit）。

vLLM与SGLang的竞争正在演变为推理引擎赛道的"Android vs iOS"：vLLM追求通用性和最大社区覆盖，SGLang追求特定场景的极致性能。两个项目的团队都来自同一所大学，私下关系紧密但技术路线日益分化。对行业而言，这种良性竞争正是推动推理成本持续下降的重要动力。

% ===========================================================
\subsection{Baseten：推理平台的融资奇观}
\label{subsec:baseten}

Baseten的融资历程本身就是一个值得写进商学院案例的故事。这家2019年成立于旧金山的推理平台公司，在\textbf{不到一年的时间内完成了三轮融资}：

\begin{itemize}[noitemsep]
  \item \textbf{2025年5月}：Series C，7500万美元，IVP领投。
  \item \textbf{2025年10月}：Series D，1.5亿美元，估值21.5亿美元，CapitalG领投。
  \item \textbf{2026年1月}：Series E，3亿美元，估值\textbf{50亿美元}，IVP和CapitalG共同领投——其中\textbf{NVIDIA单独投入1.5亿美元}。
\end{itemize}

从C轮到E轮，估值在8个月内翻了不止两倍。总融资额5.25亿美元。这种节奏在任何行业都极为罕见。

Baseten的产品定位是\textbf{推理平台}——不像Fireworks或Together那样同时提供模型API，Baseten专注于让企业在自己的基础设施上高效部署和运行模型。其核心产品是一个自动化的推理编排层：开发者上传模型，Baseten自动处理量化、批处理优化、自动扩缩容、GPU调度等所有底层细节。

客户名单的质量解释了投资人的狂热：Cursor（AI代码编辑器，估值90亿美元的现象级产品）和Notion（估值100亿美元的生产力平台）都是Baseten的核心客户。当你的推理平台在支撑着2026年最热门的AI产品时，投资人很难不掏出支票。

NVIDIA在E轮中投入1.5亿美元（占轮次的50\%）具有深层战略含义。NVIDIA正在构建一个围绕其GPU生态的"推理联盟"——通过投资Baseten、Together AI等推理平台，确保这些平台深度绑定NVIDIA硬件和软件栈（CUDA、TensorRT-LLM），从而巩固GPU在推理市场的主导地位。对Baseten而言，NVIDIA的背书不仅是资金支持，更是销售渠道的加持——NVIDIA的企业客户天然成为Baseten的潜在用户。

\begin{whybox}[为什么NVIDIA要系统性收购推理创业公司？]
NVIDIA在2024--2025年间连续收购了OctoAI（\textasciitilde\$165M, 2024年9月）和Lepton AI（"数亿美元", 2025年4月），并大额投资了Baseten、Together AI等公司。这背后的逻辑清晰而冷酷：

\textbf{控制推理软件栈}。NVIDIA的利润来源是GPU硬件，但推理软件栈决定了哪些GPU被使用、如何被使用。如果vLLM等开源引擎在AMD MI300X上跑得同样好，企业就可能转向更便宜的替代硬件。通过收购推理公司并将其技术整合到NVIDIA自有生态（NIM, TensorRT-LLM），NVIDIA确保推理优化与CUDA深度绑定。

\textbf{获取人才}。OctoAI的创始人Luis Ceze是Apache TVM（编译器框架）的核心开发者，Lepton AI的创始人Yangqing Jia是Caffe和ONNX的创始人。这些人是全球顶尖的AI系统工程师——收购公司的本质是收购大脑。

\textbf{防御AMD的追赶}。AMD的ROCm软件栈正在快速成熟，MI300X的性价比在某些推理场景下已经追平H100。NVIDIA通过掌控推理软件层来制造额外的迁移壁垒——即使硬件可以替换，被NVIDIA优化过的推理栈也很难一夜之间移植到AMD上。
\end{whybox}

% ===========================================================
\subsection{被收购与消亡：推理赛道的另一面}
\label{subsec:inference-exits}

并非每个推理创业公司都能走向IPO。事实上，这个赛道更常见的结局是被收购或关停。

\textbf{Lepton AI}的故事极其浓缩。2023年由Caffe创始人、前阿里巴巴副总裁Yangqing Jia在Cupertino创立，仅融资1100万美元，就在2025年4月被NVIDIA以"数亿美元"收购。从成立到退出不到两年，Jia的个人品牌和技术能力使得Lepton AI本质上是一个"人才收购"案例——NVIDIA买的是Jia本人及其团队，而非商业规模。

\textbf{OctoAI}走了更长的路。2019年由华盛顿大学教授Luis Ceze在西雅图创立，基于Apache TVM编译器技术构建推理优化平台，累计融资1.32亿美元。2024年9月被NVIDIA以约1.65亿美元收购——这对投资人而言几乎是平本退出（总融资1.32亿美元 vs.收购价1.65亿美元）。OctoAI的教训在于：\textbf{编译器层面的优化虽然技术含金量极高，但距离客户太远，难以独立建立商业规模}。

\textbf{Replicate}代表了又一种退出路径。2019年在Berkeley成立，累计上架超过50,000个开源模型，打造了一个"模型一键部署"的开发者平台。2025年11月，Cloudflare宣布收购Replicate（Cloudflare Press Release, 2025年11月17日）。这次收购的逻辑是互补的：Cloudflare拥有全球最大的CDN网络但缺乏AI推理能力；Replicate拥有成熟的模型部署平台但缺乏底层基础设施和企业客户渠道。

\textbf{BentoML}走的是"被生态整合"的路线。2018年成立，是最早的开源ML Serving框架之一，在Docker-style的模型打包和部署方面有深厚积累。2026年2月，BentoML被Modular收购。Modular的创始人Chris Lattner是LLVM和Swift编程语言的创始人，正在用Mojo语言构建下一代AI计算基础设施。BentoML的部署能力与Modular的编译优化能力组合，意在打造一个从模型编译到生产部署的端到端平台。

\begin{warnbox}[Banana之死：AI Wrapper的教训]
在所有推理创业故事中，Banana的结局最具警示意义。

Banana成立于2021年旧金山，定位是Serverless GPU推理平台，在Stable Diffusion爆发期吸引了超过3000个开发团队。但2024年初，创始人Erik Dunteman宣布停止运营，并在社交媒体上坦诚分享了原因：

\begin{itemize}[noitemsep]
  \item \textbf{PMF从未真正建立}：Banana服务的客户大多是"AI wrapper"——在别人模型上加一层界面的轻应用。这些客户自身的商业模式脆弱，付费意愿和能力都不稳定。
  \item \textbf{技术护城河不足}：随着vLLM、TGI等开源推理引擎成熟，Banana的核心价值——"简化GPU推理部署"——被开源工具和更大的竞争对手（Modal, RunPod, Replicate）快速复制。
  \item \textbf{创始人倦怠}：Dunteman在告别文中写道："我意识到我已经没有新的想法了，在勉强维持一个我内心不再相信的东西。"
\end{itemize}

Banana的教训是：在AI基础设施赛道，\textbf{仅靠开发者体验的差异化是不够的}——你需要要么有底层技术壁垒（如vLLM的PagedAttention），要么有规模效应（如Baseten的企业客户），要么有生态位优势（如Replicate被Cloudflare收购后获得的分发网络）。否则，当更大的玩家决定进入你的领域时，你的护城河将在几个月内消失。
\end{warnbox}

% ===========================================================
\subsection{赛道格局与前瞻}
\label{subsec:inference-landscape}

\begin{table}[htbp]
\centering
\caption{推理优化赛道创业公司全景（截至2026年3月）}
\label{tab:inference-landscape}
\small
\begin{tabularx}{\textwidth}{l c r r X}
\toprule
\rowcolor{figLightGray}
\textbf{公司} & \textbf{成立} & \textbf{估值/退出} & \textbf{总融资} & \textbf{结局} \\
\midrule
Fireworks AI      & 2022 & \$4B & \$327M & 独立; ARR$\sim$\$130M; 企业推理云 \\
Together AI       & 2022 & \$3.3B$\to$\$7.5B & \$305M+ & 独立; 洽谈\$1B融资; 全栈AI云 \\
Inferact (vLLM)   & 2023 & \$800M & \$150M & 独立; 史上最大种子轮; 开源商业化 \\
RadixArk (SGLang) & 2023 & \$400M & 未披露 & 独立; Agent推理; Accel领投 \\
Baseten           & 2019 & \$5B & \$525M & 独立; 一年三轮; NVIDIA \$150M \\
Lepton AI         & 2023 & 数亿美元 & \$11M & $\to$ NVIDIA (2025/4) \\
OctoAI            & 2019 & $\sim$\$165M & \$132M & $\to$ NVIDIA (2024/9) \\
Replicate         & 2019 & 未披露 & 未披露 & $\to$ Cloudflare (2025/11) \\
BentoML           & 2018 & 未披露 & 未披露 & $\to$ Modular (2026/2) \\
Banana            & 2021 & --- & 少量 & 关停 (2024/3) \\
\bottomrule
\end{tabularx}
\end{table}

纵观整个推理优化赛道，一个清晰的分层格局已经形成。

\textbf{第一梯队}是Fireworks AI、Together AI和Baseten——估值40亿美元以上、拥有规模化收入和优质企业客户的推理平台。它们的竞争焦点正从"谁的推理更便宜"转向"谁的平台更完整"——推理只是起点，微调、Agent编排、数据管线才是锁定客户的关键。

\textbf{第二梯队}是Inferact和RadixArk——开源推理引擎的商业化探索者。它们拥有最强的技术影响力（vLLM和SGLang合计服务了行业绝大多数推理需求），但商业模式仍在验证中。成功与否取决于一个关键问题：企业愿不愿意为开源软件的"企业版"付费？

\textbf{退出层}——Lepton AI、OctoAI、Replicate、BentoML——被大公司收购，说明推理赛道的技术人才和垂直能力具有很高的并购价值，但独立建立大规模商业体的难度极高。

\textbf{消亡层}——Banana的关停提醒我们，在基础设施赛道，没有技术护城河的"体验层"创业极其脆弱。

展望2026年下半年及更远的未来，推理赛道的演进方向将沿三条主线展开。第一，\textbf{推理与训练的融合}：随着Test-Time Compute（推理时计算）范式的成熟，推理不再是简单的"前向传播"，而是包含搜索、验证、自我修正等复杂计算——推理平台需要越来越接近训练平台的能力。第二，\textbf{硬件异构化}：NVIDIA之外，AMD MI300X、Google TPU v6、AWS Trainium3都在推理市场发力——跨硬件的推理编排将成为新的差异化方向。第三，\textbf{边缘推理的崛起}：随着Llama 3.2 1B/3B、Gemma 3等小模型的成熟，手机和笔记本端的本地推理正在从实验走向主流，这将创造一个完全不同的推理基础设施市场。

这个赛道最深层的洞察或许是：\textbf{推理优化的终极目标不是让推理更便宜，而是让推理便宜到"免费"——从而释放出一个远比今天大得多的AI应用市场}。当每一次模型调用的边际成本趋近于零，AI Agent才能真正无处不在，AI才能从一项昂贵的技术服务变成像电力一样的基础设施。推理创业公司正在竞争的，不仅是一个数十亿美元的市场，更是定义这个"AI电力"时代基础设施标准的权利。

% ===========================================================
\subsection{小众推理服务：长尾生态}
\label{subsec:inference-longtail}
% ===========================================================

在Fireworks、Together、Baseten三强之外，一批更小的推理API提供商正在覆盖不同的价格点和用户群。它们共同构成了推理赛道的长尾生态。

\subsubsection{DeepInfra}
DeepInfra（2022 年，Palo Alto）是一家精简高效的推理API平台，团队仅约10人，却服务着大量中小开发者。2025年完成\$1800万Series A融资，随后又完成\$5400万Series B，累计融资约\$7200万。DeepInfra以\textbf{极致性价比}著称——在OpenRouter等推理聚合平台的比价中，DeepInfra的定价通常是最低之列，DeepSeek-R1等热门模型的推理成本比官方API低30--50\%。DeepInfra的策略很清晰：不做平台，不做微调，只做一件事——\textbf{把开源模型的推理做到最便宜}。

\subsubsection{Novita AI}
Novita AI（2023 年，Singapore）提供覆盖LLM、图像生成、语音、视频的\textbf{多模态推理API}，号称接入超过100个模型和10,000个图像模型。2025年与SGLang达成战略合作，采用其推理引擎优化后端性能。团队约10人，年营收约\$110万。Novita AI的差异化在于其多模态广度——当大多数推理平台聚焦文本LLM时，Novita同时覆盖了Stable Diffusion、语音合成等长尾AI应用场景，服务于图像/视频生成领域的开发者。

\subsubsection{SiliconFlow（硅基流动）}
SiliconFlow（2023 年，Singapore/中国）是中国推理加速赛道的代表性创业公司。自研推理加速引擎支持NVIDIA H100/H200和AMD MI系列，提供Serverless端点、专用GPU实例和BYOC（自带集群）三种部署模式，API兼容OpenAI格式。2025年7月完成\$2080万Series A融资。SiliconFlow的定位是\textbf{高性能推理基础设施}——不是简单的API转售，而是在底层做算子优化和推理调度，目标客户是需要大规模生产级推理的中国和全球AI企业。

\subsubsection{芯片厂商的推理云服务}
一个值得关注的趋势是AI芯片创业公司纷纷推出自营推理API，从``卖芯片''转向``卖推理''。\textbf{GroqCloud}在被NVIDIA收购前已积累超过250万开发者，以LPU的确定性低延迟推理著称；\textbf{Cerebras Inference}基于WSE-3芯片提供云端推理，在Llama 3.1 405B上实现超过1000 tok/s的吞吐，是GPU方案的数倍；\textbf{SambaNova Cloud}基于SN50 RDU芯片提供推理API，在agentic AI场景下号称比B200快5倍。这些芯片厂商的推理服务本质上是\textbf{用自研硅片的性能优势补贴推理定价}，以获取开发者生态和商业验证。对于用户而言，它们提供了GPU之外的另一个选择；对于芯片厂商自身，推理云是最直接的商业化路径——比卖硬件更快地产生收入和用户反馈。


% §2.5 AI Factories and Supercompute Clusters
% ch02-s5-factories.tex  — §2.5 AI Factories and Supercompute Clusters
% Part of Chapter 2: Infrastructure Layer
% Input from: ch02-infrastructure.tex (or main file)

\section{AI 工厂与超级集群}
\label{sec:infra-factories}

NVIDIA CEO Jensen Huang 在 GTC 2025 上提出了一个定义时代的比喻：
\textbf{``上一次工业革命的原材料是水，产品是电；
这一次，原材料是数据和电力，产品是 token。''}
到了 GTC 2026，他进一步将这一概念具象化——
``AI 工厂''（AI Factory）不再是隐喻，而是一种全新的工业设施类型：
吞入海量数据和百兆瓦级电力，
产出的不是钢铁或电能，而是智能——以 token 为单位计量的数字商品。

这一概念的背后是一组令人瞠目的数字。
2025 年，全球五大超算厂商（Microsoft、Google、Meta、Amazon、Oracle）
的资本开支合计约 \$3,810 亿；
进入 2026 年，这一数字预计飙升至 \$6,350--6,650 亿，
同比增长 67--74\%。
Moody's 预测超算厂商 capex 正逼近 \$7,000 亿大关。
Dell'Oro Group 的数据显示，
2025 年全球数据中心资本开支已同比猛增 57\%，
且 2026 年的增长势头丝毫未减。
AI 算力的建设规模，
已经可以与 20 世纪初的铁路和电力基础设施投资相提并论。

\begin{keybox}[全球超算厂商资本开支演进（2022--2026E）]
\small
\renewcommand{\arraystretch}{1.3}
\begin{tabularx}{\linewidth}{lXXXXX}
\toprule
\rowcolor{figLightGray}
\textbf{公司} & \textbf{2022} & \textbf{2023} & \textbf{2024} & \textbf{2025} & \textbf{2026E} \\
\midrule
Microsoft   & \$24B  & \$32B  & \$56B  & \$80B   & \$110--120B \\
Google      & \$31B  & \$32B  & \$52B  & \$75B   & \$100--110B \\
Meta        & \$32B  & \$28B  & \$39B  & \$66--72B & \$95--115B \\
Amazon      & \$58B  & \$48B  & \$74B  & \$100B  & \$130--150B \\
Oracle      & \$7B   & \$8B   & \$12B  & \$30B+  & \$50--60B \\
\midrule
\textbf{合计} & \$152B & \$148B & \$233B & \$381B & \$635--665B \\
\bottomrule
\end{tabularx}

\vspace{4pt}
\noindent\textit{数据来源：各公司财报、Futurum Research、Yahoo Finance 汇编。
2026E 为分析师共识估计。}
\end{keybox}

以下逐一剖析全球最具代表性的 AI 超级集群项目。

% =====================================================================
\subsection{xAI Colossus：122 天建成的超级计算机}
\label{subsec:xai-colossus}

2024 年 7 月，Elon Musk 在 Memphis, Tennessee 启动了一场史无前例的``闪电战''：
从空地到 10 万块 H100 GPU 上线，仅用了 \textbf{122 天}。
这就是 xAI 的 Colossus 超算集群——
以速度而非规模率先震撼了业界。

Colossus 的第一阶段部署了 10 万块 H100；
到 2024 年底升级为 15 万块 H100 加 5 万块 H200，
并引入了 3 万块 NVIDIA GB200。
值得注意的是，
Colossus 选择了 \textbf{NVIDIA Spectrum-X 以太网}而非传统的 InfiniBand 互联——
这是一个重大的技术赌注。
Spectrum-X 基于 RoCE（RDMA over Converged Ethernet）协议，
理论上可以在降低成本的同时实现更灵活的网络拓扑，
但在超大规模训练中是否能匹敌 InfiniBand 的可靠性，
至今仍是行业争论的焦点。

到 2025 年底，xAI 已购入第三栋建筑，
目标是将总功率扩展至 2\,GW，
GPU 数量推向 100 万块。
2026 年 2 月的最新消息显示，
Memphis 园区的 GPU 总量已接近 55.5 万块。
2026 年 3 月，xAI 又提交了一份 \$6.59 亿的扩建许可，
为 Colossus 2 新增一栋 31.2 万平方英尺的四层建筑。
Musk 的目标清晰且激进：
\textbf{在 AGI 竞赛中，算力就是弹药，速度就是战略}。

% =====================================================================
\subsection{Stargate：史上最大 AI 基础设施计划}
\label{subsec:stargate}

2025 年 1 月，OpenAI、SoftBank 和 Oracle 联合宣布了 \textbf{Stargate 计划}——
一项总投资 \$5,000 亿、为期四年的 AI 基础设施建设项目，
被多家媒体冠以``人类历史上最大规模的 AI 基础设施投资''。

Stargate 的旗舰站点位于 Texas 州 Abilene，
规划用电量达 1.2\,GW，
将部署 40 万块 NVIDIA GB200 GPU，
运行在 Oracle OCI Zettascale10 架构之上，
峰值算力目标为 \textbf{16 zettaFLOPS}——
这一数字是 2024 年全球最强超算 Frontier 的数千倍。

到 2025 年 9 月，OpenAI、Oracle 和 SoftBank 宣布新增五个站点，
称 Stargate 的建设``进度超前''，
有望在 2025 年底前锁定全部 \$5,000 亿投资和 10\,GW 总算力承诺。
然而，2026 年 2 月传出的消息给这幅宏伟蓝图蒙上了阴影：
三方伙伴之间就\textbf{数据中心的最终控制权}产生了分歧——
OpenAI、Oracle 和 SoftBank 各有不同的战略诉求，
项目部分环节出现延迟。
这一插曲揭示了超大型 AI 基础设施项目的深层挑战：
技术和资金固然重要，
但\textbf{治理结构与利益对齐}才是决定项目成败的关键变量。

% =====================================================================
\subsection{CoreWeave：从加密矿场到 AI 云基础设施巨头}
\label{subsec:coreweave-cluster}

CoreWeave 的故事是 AI 时代最不可思议的转型叙事之一。
这家公司从加密货币挖矿起步，
转型为 NVIDIA GPU 云服务商后一路狂飙，
2025 年 3 月在 IPO 前夕与 OpenAI 签下 \$11.9B 合同，
随后 9 月又追加 \$6.5B，合计 \textbf{\$22.4B} 的 OpenAI 算力承诺。
截至 2025 年 Q3，CoreWeave 的营收积压（backlog）已翻倍至 \$55.6B，
公司将 2025 年全年 capex 设定为 \$23B。

CoreWeave 的战略不止于美国本土。
在英国，Crawley 和 Docklands 的数据中心已投入运营，
目标是到 2030 年将全球算力规模扩展至 \textbf{5+\,GW}。
CoreWeave 的核心优势在于\textbf{极致的 GPU 密度和裸金属灵活性}——
不同于 AWS 或 Azure 的通用云，
CoreWeave 从底层就为 AI 工作负载优化，
这让它在训练和推理的性价比上形成了差异化竞争力。
其风险也同样突出：
客户集中度极高（OpenAI 占据绝对多数），
一旦 OpenAI 转向自建算力或分散供应商，
CoreWeave 的营收基础将面临严峻考验。

% =====================================================================
\subsection{Meta RSC 及扩展：开放计算的超级工厂}
\label{subsec:meta-rsc}

Meta 在 AI 基础设施上的投入力度，
足以让多数独立国家的科技预算黯然失色。
2025 年，Meta 宣布全年 capex 为 \$66--72B，
2026 年预计进一步攀升至近 \$100B。

Meta 的超级集群布局分为两条线：
\textbf{Prometheus} 和 \textbf{Hyperion}。
Prometheus 是 Meta 首个 GW 级 AI 集群，
计划 2026 年上线，采用了一种名为
\textbf{Backend Aggregation (BAG)} 的创新架构——
通过 Disaggregated Schedule Fabric (DSF) 和
Non-Scheduled Fabric (NSF) 两种网络结构的融合，
实现 GW 级集群的高效互联。
2026 年 2 月，Meta 工程博客详细披露了 Prometheus 的技术细节，
标志着这一项目从概念走向实施。

Hyperion 则是更远期的愿景：
一个 \textbf{5\,GW} 的超级数据中心园区，
计划部署``数百万''块 NVIDIA Blackwell 和下一代 Rubin GPU。
若 Hyperion 全面建成，
其单一站点的用电量将超过许多小型国家的全国用电。

Meta 的独特之处在于其\textbf{开放计算理念}（Open Compute Project）。
不同于 Google 和 Amazon 的自研+封闭路线，
Meta 将数据中心设计开源，
推动了整个行业的服务器和冷却技术进步。
这种``开放''策略既是技术哲学，
也是降低供应链风险的务实选择——
当你公开设计标准时，
供应商之间的竞争自然会压低成本。

% =====================================================================
\subsection{Google TPU Pods：垂直整合的算力之路}
\label{subsec:google-tpu}

在所有超大规模玩家中，
Google 走了一条最彻底的\textbf{垂直整合}路线：
自研芯片（TPU）、自研网络（Jupiter）、
自研编译器（XLA）、自研框架（JAX），
端到端控制从硅片到模型的每一层。

TPU v5p 每个 Pod 包含 8,960 颗芯片，
单芯片峰值算力 459\,TFLOPS，
通过 Google 的定制 ICI（Inter-Chip Interconnect）网络互联。
2025 年发布的 \textbf{TPU v6e Trillium} 实现了单芯片 \textbf{4.7 倍}的性能提升，
在密集 LLM 训练中约快 4 倍，
推理吞吐提升 2.8--2.9 倍。
每颗 Trillium 配备 32\,GB HBM，内存带宽约 1,600\,GB/s，
单个 Pod 可扩展至 256 颗芯片。

Anthropic 签下了 Google 历史上最大的 TPU 合同——
数十万块 Trillium 芯片，计划到 2027 年扩展至 100 万块。
Midjourney 从 GPU 迁移至 TPU 后推理成本降低了 65\%。
这些数据表明，
TPU 正在从 Google 的``内部工具''演变为
\textbf{具有外部竞争力的算力平台}。

Google 还提出了 \textbf{AI Hypercomputer} 的整合架构概念，
将 TPU Pod、GPU 节点、存储和网络统一编排，
为不同规模和类型的 AI 工作负载提供最优资源分配。
2026 年发布的 \textbf{Ironwood（TPU v7）}
则进一步推高了 Google 自研芯片的性能天花板。

% =====================================================================
\subsection{Oracle OCI：Stargate 背后的基础设施骨干}
\label{subsec:oracle-oci}

Oracle 在 AI 基础设施竞赛中的崛起是最出人意料的故事之一。
曾经被视为``传统企业软件公司''的 Oracle，
凭借 OCI（Oracle Cloud Infrastructure）在 AI 云市场中杀出了一条血路。

2025 年 10 月，Oracle 发布了 \textbf{Zettascale10} 集群——
号称``云上最大的 AI 超级计算机''，
包含 131,072 块 NVIDIA Blackwell GPU，
峰值性能达 10 倍于前代 Zettascale 集群的 zettaFLOPS 级算力。
更重要的是，Oracle 推出了自研的 \textbf{Acceleron RoCE 网络架构}，
这是一种基于以太网的高速互联方案，
直接挑战 InfiniBand 在超大规模训练中的垄断地位。
Oracle 声称 Zettascale10 可扩展至 \textbf{80 万块 GPU} 的超级集群。

Oracle 的 capex 增长同样惊人：
2025 年的资本支出预计是前一年的三倍，
反映了 Stargate 项目和其他 AI 客户需求的强力拉动。
作为 Stargate 的核心基础设施提供商，
Oracle 已将自身从数据库公司重新定位为
\textbf{AI 时代的算力基础设施供应商}——
这一转型的成功与否，将决定 Oracle 未来十年的市场地位。

% =====================================================================
\subsection{ByteDance/火山引擎：中国 AI 算力的隐形巨头}
\label{subsec:bytedance-infra}

在中国 AI 云市场中，
ByteDance 旗下的火山引擎（Volcengine）悄然成长为仅次于阿里云的第二大玩家。
其 AI 应用豆包（Doubao）的月活跃用户已达 \textbf{1.7 亿}，
庞大的用户规模反向驱动了对算力基础设施的饥渴需求。

更引人注目的是 ByteDance 的\textbf{海外算力布局}。
2026 年 3 月，多家媒体披露 ByteDance 正在
马来西亚等海外市场悄悄建设大规模 NVIDIA GPU 集群，
投资规模约 \$25 亿，
采购的是 NVIDIA 最新的 Blackwell 系列芯片。
这一布局的逻辑清晰：
在中美芯片管制的大背景下，
\textbf{海外节点既能规避出口管制限制，
又能服务于 TikTok 等全球化产品的 AI 推理需求}。

ByteDance 的 AI 基础设施战略体现了中国科技公司的一种普遍模式：
在国内依靠自研芯片和国产替代方案应对制裁压力，
在海外则尽可能地获取最先进的 NVIDIA 硬件。
这种``双轨制''将是理解未来全球 AI 算力格局的重要视角。

% =====================================================================
\subsection{中东 AI 数据中心：第三极的崛起}
\label{subsec:middle-east-dc}

\begin{whybox}[为什么中东成为 AI 算力第三极？]
三个结构性因素推动中东成为继美国和中国之后的全球 AI 算力第三极：
\begin{itemize}[noitemsep]
  \item \textbf{主权财富基金的战略转向}——
        沙特 PIF、阿联酋 Mubadala 等主权基金将 AI 基础设施
        视为``后石油时代''经济转型的核心支柱，
        资金规模和决策速度远超市场化融资。
  \item \textbf{能源禀赋的先天优势}——
        中东拥有全球最低的发电成本（天然气+光伏），
        可以为 GW 级数据中心提供极具竞争力的电力价格。
  \item \textbf{地缘政治的``中间人''角色}——
        中东国家同时与中美保持良好关系，
        可以成为全球 AI 模型训练和推理的``中立枢纽''，
        吸引双方都无法直接服务的市场。
\end{itemize}
\end{whybox}

沙特阿拉伯的布局最为激进。
在暂停了耗资 \$500 亿的 NEOM ``The Line'' 项目后，
沙特将资源迅速转向了 AI 基础设施。
国有 AI 公司 \textbf{HUMAIN}
（由公共投资基金 PIF 控股）正在全国规划 211 个站点，
建设 GW 级 AI 数据中心网络。
2026 年 3 月，HUMAIN 与 MIS 签署了 \$5 亿的数据中心设计建造合同，
同时与 Qualcomm 合作部署 AI200 系列推理芯片。
NEOM 本身仍保留了 \$50 亿、1.5\,GW 的数据中心规划（目标 2028 年）。
沙特政府已将 2026 年定为``AI 之年''。

阿联酋方面，G42（背靠 Abu Dhabi 的 AI 巨头）
与 Microsoft 达成了 \$15.2B 的战略合作，
涵盖 Azure 云部署和 AI 模型训练。
G42 正在中东和中亚地区建设多个数据中心集群，
目标是成为``非美非中''世界的首选 AI 算力提供商。

% =====================================================================

\begin{corebox}[全球 AI 超算竞赛：军备竞赛式投入]
2026 年的全球 AI 基础设施竞赛已具备``军备竞赛''的所有特征：
\begin{itemize}[noitemsep]
  \item \textbf{规模指数级增长}——
        从 2022 年的 \$1,520 亿到 2026 年的 \$6,500 亿以上，
        四年翻了 4 倍多。
  \item \textbf{国家意志介入}——
        美国 Stargate 获白宫背书，
        沙特以主权基金推动，
        中国以``新基建''政策牵引。
  \item \textbf{赢家通吃的逻辑}——
        算力优势直接转化为模型能力优势，
        而模型能力优势又吸引更多数据和用户，
        形成``飞轮效应''。
  \item \textbf{退出成本极高}——
        一旦开始建设 GW 级数据中心，
        停止投入意味着沉没成本和竞争力的双重损失。
\end{itemize}
这不是一场可以选择观望的竞赛。
正如 Jensen Huang 在 GTC 2026 上所言：
``世界正在从通用计算时代进入 token 经济时代——
每一家公司都需要自己的 AI 工厂。''
\end{corebox}

% =====================================================================

\begin{warnbox}[能源危机：AI 的电力饥渴]
AI 超级集群的狂飙突进正将全球能源系统推向极限：

\begin{itemize}[noitemsep]
  \item \textbf{电力消耗}——
        美国数据中心 2024 年耗电 183\,TWh，
        占全国总用电量约 4.4\%。
        到 2030 年，这一比例预计翻倍至 8--9\%。
  \item \textbf{核能回归}——
        Microsoft 以 \$16B、20 年合约推动三里岛核电站重启（835\,MW，目标 2028 年）；
        Google 与 Kairos Power 签署首个企业级 SMR 合同（500\,MW，2030+ 上线）；
        Amazon 投资超过 \$20B 将 Susquehanna 核电站转型为 AI 园区。
        科技巨头的核电合同总额已超过 \$52B。
  \item \textbf{碳排放悖论}——
        科技公司一方面承诺碳中和，
        另一方面却在建设数百万千瓦的化石燃料+核能发电设施。
        Google 2023 年碳排放同比增长 13\%，
        Microsoft 增长 29\%——
        而 AI 基础设施的大规模建设才刚刚开始。
  \item \textbf{水资源压力}——
        液冷和蒸发冷却技术虽然提升了 PUE，
        但 GW 级数据中心每年消耗数十亿升水，
        在水资源紧张的地区（如美国西部、中东）引发日益激烈的社区反对。
\end{itemize}

AI 工厂的``产品是 token''，
但其``副产品''——碳排放、热污染和水资源消耗——
正在成为这个时代最大的外部性之一。
\end{warnbox}

从 Memphis 的闪电施工到 Abilene 的万亿投资，
从 Crawley 的 GPU 密集机房到利雅得的主权 AI 工厂，
2022--2026 年的 AI 基础设施建设已经重塑了全球的能源版图、
供应链格局和地缘政治博弈。
Jensen Huang 的``AI 工厂''比喻并不夸张——
这些设施正在成为 21 世纪的发电站，
只不过它们输出的不是电力，而是智能。

下一节将讨论连接这些超级集群的高速网络和互联技术
（\S\ref{sec:infra-networking}），
没有它们，再多的 GPU 也只是一堆孤立的硅片。


% §2.6 Networking and Interconnect + §2.7 CDN and Edge AI
% ============================================================
% §2.6 Networking and Interconnect + §2.7 CDN and Edge AI
% File: ch02-s67-network-edge.tex
% ============================================================

\section{网络与互联}
\label{sec:infra-networking}

数据的高速公路：为什么互联是AI基础设施的隐形瓶颈？当一个训练集群扩展到数万块GPU时，
单块芯片的算力已不是第一约束——GPU之间``搬运数据''的速度才是。一个100,000 GPU集群
每次all-reduce同步可产生PB级流量；若网络带宽不足或延迟过高，GPU将大量时间空转等待。
这正是为什么2024--2026年间，光互联、高速网络fabric成为风险投资的新宠赛道。

\subsection{光互联三巨头}

传统电互联（铜缆）在400~Gbps以上面临严重的功耗墙和距离限制。共封装光学
（Co-Packaged Optics, CPO）将光收发器直接集成到芯片封装中，从根本上突破了
``silicon beachfront''约束——即芯片边缘可用于I/O的物理面积有限这一问题。
三家创业公司在这一赛道形成了鲜明的技术路线分化。

\begin{corebox}{光互联三巨头对比：Lightmatter vs Ayar Labs vs Celestial AI}
\begin{tabularx}{\textwidth}{l X X X}
\toprule
\rowcolor{figLightGray}
\textbf{维度} & \textbf{Lightmatter} & \textbf{Ayar Labs} & \textbf{Celestial AI} \\
\midrule
成立 & 2017, Mountain View & 2015, San Jose & Santa Clara \\
创始人背景 & Nick Harris, MIT光子学PhD & MIT/Berkeley/Colorado光子学联合团队 & 前Luxtera高管 \\
核心产品 & Passage CPO芯片 & TeraPHY光I/O chiplet & Photonic Fabric \\
技术路线 & DWDM多波长光引擎 & 硅光chiplet直接共封装 & 光子交换+远距光I/O \\
最新里程碑 & 1.6 Tbps/光纤（世界纪录）; Passage L20 6.4 Tbps双向光引擎 & Series E \$500M @ \$3.75B (2026/3); 量产CPO chiplet & 2025/12 被Marvell \$3.25B收购 \\
融资总额 & $\sim$\$850M (Series D \$400M @ \$4.4B) & \$870M & Series C1 \$250M @ \$2.5B \\
投资者亮点 & Google, Viking Global & AMD, Intel, NVIDIA三巨头均投 & 未披露 \\
\bottomrule
\end{tabularx}
\end{corebox}

\textbf{Lightmatter}是光互联领域势头最猛的玩家。2026年3月，公司宣布Passage
L20光引擎——单芯片双向带宽达6.4~Tbps，同时与Qualcomm合作实现了单根光纤
1.6~Tbps的世界纪录（16波长DWDM技术，较传统FR-4方案实现8倍带宽密度提升）。
更值得注意的是，Lightmatter在OCP（Open Compute Project）内牵头发起了CPO开放
参考架构计划，意图将光互联从专有技术推向行业标准——这一策略若成功，将使其成为
光互联领域的``ARM''。

\textbf{Ayar Labs}走的是硅光chiplet共封装路线。2026年3月，公司完成\$500M的
Series E融资，估值\$3.75B。其独特优势在于同时获得了AMD、Intel和NVIDIA三大芯片
巨头的投资——这在半导体行业极为罕见，意味着TeraPHY chiplet有望成为跨平台的
光I/O标准组件。NVIDIA在Ayar Labs融资的前一天还宣布向Coherent和Lumentum注资
\$4B，显示出整个光互联供应链正在获得前所未有的资本投入。

\textbf{Celestial AI}的Photonic Fabric采用了最激进的架构：不仅实现芯片间光互联，
还提供光子交换和远距内存池化（论文展示了32~TB共享HBM3E内存 + 115~Tbps全对全
数字交换的能力）。然而，公司在2025年12月被Marvell以\$3.25B收购，成为光互联领域
首个重大并购退出案例。这一结局说明：即使技术领先，创业公司在面对半导体巨头的
整合需求时，被收购可能比独立上市更现实。

\subsection{网络芯片与fabric创业}

\textbf{Enfabrica}（2019年，Mountain View）开发了ACF-S ``Millennium'' SuperNIC，
单芯片3.2~Tbps带宽，定位为GPU集群的以太网fabric互联。公司累计融资\$290M，
但在2025年9月被NVIDIA以超\$900M的价格``acquihire''——NVIDIA获得了CEO Rochan
Sankar及核心团队和技术授权，但并非传统意义上的全资收购。这笔交易反映了NVIDIA
在以太网AI网络领域加速布局的战略意图。

\textbf{Astera Labs}（Santa Clara）是AI网络领域少有的IPO成功案例。2024年3月
在NASDAQ上市（代码ALAB），首日股价暴涨72\%，成为2024年首个大型科技IPO。公司
专注于PCIe/CXL智能连接芯片，融资\$713M，上市时市值\$9.46B。Astera Labs的成功
说明：在AI基础设施的``管道工''层（连接芯片、retimer、信号调理），同样存在巨大的
价值捕获机会。

\subsection{InfiniBand vs 以太网：一场正在发生的大转折}

\begin{keybox}{InfiniBand → 以太网：关键转折数据}
\begin{itemize}
  \item \textbf{2023年}：InfiniBand占据AI后端网络$\sim$80\%份额，几乎是唯一选择
  \item \textbf{2025年中}：以太网首次超越InfiniBand（Dell'Oro Group数据）
  \item \textbf{2025全年}：以太网AI后端交换机销售额\textbf{超过InfiniBand 2倍以上}，
    占AI集群数据中心交换机总销售的三分之二以上
  \item \textbf{以太网交换机销售额}在2025年Q4同比增长超过3倍
  \item \textbf{UEC Spec 1.0}：2025年6月发布，标准化高性能AI以太网
  \item \textbf{领军厂商}：NVIDIA和Celestica合计占以太网AI交换机近50\%份额
\end{itemize}
\end{keybox}

这场转折的速度令人惊讶。仅仅两年前，InfiniBand还被视为大规模AI训练的唯一可行
网络方案——NVIDIA的DGX SuperPOD默认使用InfiniBand，Meta最初也用IB训练LLaMA。
但到2025年底，局势已彻底逆转。

\begin{whybox}{为什么以太网正在赢得AI网络之战？}
三个结构性因素驱动了这一转变：

\textbf{1. 超大规模云厂商的集体选择。}Amazon、Microsoft、Meta、Oracle、xAI
全部选择以太网作为下一代AI集群的后端网络。当你的五大客户都选择同一个方向时，
市场重心的转移就是时间问题。这些厂商的核心诉求是：多供应商竞争（避免被NVIDIA
锁定）、与现有数据中心网络架构统一、以及更低的总体拥有成本（TCO）。

\textbf{2. 技术差距正在快速缩小。}Ultra Ethernet Consortium（UEC）在2025年6月
发布了Spec 1.0，系统性地解决了以太网在AI训练场景中的短板：包括无损传输、
精确拥塞控制、以及集合通信原语的硬件卸载。Arista CEO Jayshree Ullal明确表示：
``以太网已经将InfiniBand推到一边。''

\textbf{3. 流量模式从训练向推理迁移。}随着AI工作负载从集中式训练转向分布式
推理，网络流量模式发生了根本变化——推理场景更适合以太网的弹性扩展模型，
而非InfiniBand的紧耦合拓扑。
\end{whybox}

InfiniBand并未消亡——在追求极致延迟的小规模训练集群中（如部分CoreWeave部署
和NVIDIA DGX Cloud），它仍有不可替代的优势。但在10万GPU以上的超大集群中，
以太网已成为事实标准。

%% ============================================================

\section{CDN 与边缘 AI}
\label{sec:infra-edge}

如果说上一节讨论的是数据中心\emph{内部}的高速公路，本节关注的则是将AI推理能力
\emph{推向用户身边}的最后一公里。边缘AI的核心命题是：当模型足够小、硬件足够强时，
推理不必发生在远方的数据中心——它可以在CDN边缘节点、手机芯片、甚至笔记本电脑上
完成。

\subsection{CDN边缘推理}

\textbf{Cloudflare Workers AI}是边缘推理的典型代表。2023年9月发布以来，
Cloudflare将AI推理能力部署到全球300+边缘节点，提供真正的serverless推理服务——
开发者无需管理GPU、无需容量规划，一次API调用即可在离用户最近的节点运行模型。
截至2026年3月，平台已支持100+模型（包括Moonshot Kimi K2.5等大模型），覆盖
LLM、图像生成、语音识别、embedding等全品类。Cloudflare还收购了Replicate，
进一步补全了模型托管能力。Workers AI的核心优势在于：与Cloudflare既有的CDN、
Workers、D1数据库、R2存储无缝集成，形成了``推理即基础设施''的完整开发者体验。

\textbf{Fastly}从传统CDN转型边缘AI。其Compute@Edge平台新增了AI Gateway功能，
支持多AI提供商路由——开发者可以在Fastly边缘节点上统一调用不同模型API，
实现负载均衡、故障切换和成本优化。2026年2月，Fastly股价单日上涨42\%，
市场对其AI边缘战略投出了信任票。

\textbf{Vercel AI SDK}虽非严格意义上的边缘推理平台，却已成为TypeScript/JavaScript
生态中构建LLM应用的事实标准。2026年3月，npm周下载量达$\sim$350万次；v0 AI
辅助开发平台拥有350万用户。Vercel的战略定位是``AI应用的部署层''——通过Edge
Runtime + AI SDK + v0，覆盖了从模型调用到前端渲染的全链路。

\subsection{端侧AI芯片}

\textbf{Qualcomm Cloud AI 200}专注于推理场景，配备768~GB LPDDR5X内存，
定位数据中心级推理加速（而非训练）。其最引人注目的合作是与沙特HUMAIN签署的
200~MW数据中心推理部署协议——这标志着AI推理基础设施正在成为主权国家级别的
战略投资。Qualcomm在边缘推理领域的差异化优势在于其移动芯片的能效比积累：
从手机到数据中心，同一架构的功耗优势具有可迁移性。

\textbf{Apple Neural Engine}在M系列芯片中持续进化。M5芯片的架构突破在于：
每个GPU核心内置Neural Accelerator，实现了4倍AI性能提升（相比M4）。Apple的
边缘AI战略有两个独特之处：一是Apple Intelligence将大模型能力直接嵌入操作系统
层面（而非App层），二是Private Cloud Compute（PCC）架构——当端侧算力不足时，
请求被路由到Apple自有数据中心的Apple Silicon服务器上处理，全程端到端加密，
甚至Apple自己也无法访问用户数据。这种``端云一体+隐私优先''的模式是Google和
Microsoft都难以复制的。

\begin{warnbox}{Intel Gaudi的教训：价格优势不够}
Intel Gaudi 3（源自2019年收购的Habana Labs）定价仅\$15K，是H100的一半。
然而，2025年出货目标下调30\%，全年营收约\$500M，远不及NVIDIA单季度的AI收入。

\textbf{教训何在？}在AI加速器市场，\emph{价格低}不等于\emph{赢}。NVIDIA的护城河
不在于硬件本身，而在于CUDA生态——数百万开发者、数千个优化库、以及十年积累的
软件栈。Gaudi 3的硬件性能尚可，但软件生态（包括编译器成熟度、框架兼容性、
社区支持）与CUDA生态存在代差。这对所有AI芯片挑战者（包括Cerebras、Groq、
Tenstorrent等）都是同样的警示：\textbf{芯片是入场券，生态才是护城河。}
\end{warnbox}

\subsection{边缘AI的格局判断}

边缘AI正在形成三个层次的分化：\textbf{CDN边缘}（Cloudflare、Fastly）负责轻量
推理和API路由，延迟优势明显但模型尺寸受限；\textbf{设备端}（Apple、Qualcomm）
负责隐私敏感和实时性要求极高的场景，受限于芯片面积和功耗但用户体验最优；
\textbf{近边缘数据中心}（Qualcomm Cloud AI、区域推理集群）填补中间地带，
适合中等规模模型的低延迟推理。2026年的共识是：没有单一方案能通吃所有场景，
``端-边-云''协同推理将成为主流架构。


% §2.8 Ecosystem Analysis — to be written after cross-section review
\section{生态位分析与竞争格局}
\label{sec:infra-landscape}

% Infrastructure layer competitive landscape market map
\begin{figure}[htbp]
\centering
\begin{tikzpicture}[
  every node/.style={font=\sffamily\scriptsize},
  bubble/.style={draw=none, fill opacity=0.7, text opacity=1},
]
% --- Axes ---
\draw[-{Stealth[length=3mm]}, thick, figGray]
  (0,0) -- (13.5,0) node[right, font=\sffamily\small] {收入规模};
\draw[-{Stealth[length=3mm]}, thick, figGray]
  (0,0) -- (0,11.5) node[above, font=\sffamily\small] {技术层级};

% X-axis tick labels (revenue)
\foreach \x/\lab in {
  1.5/{\$0},
  4/{\$100M},
  6.5/{\$500M},
  9/{\$1B},
  11.5/{\$5B+}}{
  \node[below, figGray, font=\sffamily\tiny] at (\x, -0.1) {\lab};
  \draw[figGray!40] (\x, 0) -- (\x, 0.1);
}

% Y-axis layer labels (technology tiers)
\foreach \y/\lab in {
  1.8/{Edge \& CDN},
  3.8/{网络互连},
  5.8/{推理优化},
  7.8/{训练基础设施},
  9.8/{芯片},
  11.2/{GPU 云}}{
  \node[left, figGray, font=\sffamily\tiny, text width=1.8cm, align=right] at (-0.1, \y) {\lab};
  \draw[figGray!15, dashed] (0, \y) -- (13, \y);
}

% --- Legend ---
\begin{scope}[shift={(9.2, 10.8)}]
  \fill[figBlue, opacity=0.7]   (0, 0) circle (3pt);
  \node[right, font=\sffamily\tiny] at (0.2, 0) {独立};
  \fill[figRed, opacity=0.7]    (1.8, 0) circle (3pt);
  \node[right, font=\sffamily\tiny] at (2.0, 0) {被收购};
  \fill[figGreen, opacity=0.7]  (3.4, 0) circle (3pt);
  \node[right, font=\sffamily\tiny] at (3.6, 0) {IPO/上市};
  \fill[figGray, opacity=0.5]   (0, -0.5) circle (3pt);
  \node[right, font=\sffamily\tiny] at (0.2, -0.5) {关停/转型};
  \node[right, font=\sffamily\tiny] at (2.0, -0.5) {气泡大小 $\propto$ 估值};
\end{scope}

% ====== GPU Cloud (y ~ 11.2) ======
% CoreWeave — $35B valuation, ~$2B rev → x≈10, big bubble
\fill[figBlue, bubble] (10, 11.2) circle (14pt);
\node[above right, figBlue!80!black, font=\sffamily\tiny\bfseries] at (10.5, 11.4) {CoreWeave};
% Lambda Labs — $2.8B val, ~$300M rev
\fill[figBlue, bubble] (5.2, 11.0) circle (8pt);
\node[below, font=\sffamily\tiny] at (5.2, 10.5) {Lambda};
% Together AI — $3.3B val, ~$100M rev
\fill[figBlue, bubble] (4, 11.4) circle (8pt);
\node[above, font=\sffamily\tiny] at (4, 11.9) {Together};
% Crusoe Energy — $3.3B val
\fill[figBlue, bubble] (3.5, 10.8) circle (8pt);
\node[below left, font=\sffamily\tiny] at (3.2, 10.5) {Crusoe};
% Voltage Park — smaller
\fill[figBlue, bubble] (2.2, 11.5) circle (5pt);
\node[above, font=\sffamily\tiny] at (2.2, 11.8) {Voltage Park};
% Nebius — IPO
\fill[figGreen, bubble] (6, 11.5) circle (7pt);
\node[above, font=\sffamily\tiny] at (6, 11.9) {Nebius};

% ====== Chips (y ~ 9.8) ======
% Cerebras — IPO, $8B+
\fill[figGreen, bubble] (3.5, 9.8) circle (11pt);
\node[above, font=\sffamily\tiny\bfseries] at (3.5, 10.3) {Cerebras};
% Groq — $2.8B val
\fill[figBlue, bubble] (4.5, 9.5) circle (8pt);
\node[below right, font=\sffamily\tiny] at (4.8, 9.2) {Groq};
% SambaNova — $5.1B val
\fill[figBlue, bubble] (5.8, 10.0) circle (9pt);
\node[above right, font=\sffamily\tiny] at (6.2, 10.2) {SambaNova};
% d-Matrix
\fill[figBlue, bubble] (1.8, 9.6) circle (5pt);
\node[below, font=\sffamily\tiny] at (1.8, 9.2) {d-Matrix};
% Tenstorrent — $2.8B
\fill[figBlue, bubble] (3, 9.2) circle (7pt);
\node[below, font=\sffamily\tiny] at (3, 8.7) {Tenstorrent};
% Etched
\fill[figBlue, bubble] (2.0, 10.2) circle (5pt);
\node[left, font=\sffamily\tiny] at (1.6, 10.2) {Etched};
% 寒武纪 (Cambricon) — IPO in China
\fill[figGreen, bubble] (7.5, 9.8) circle (9pt);
\node[above, font=\sffamily\tiny] at (7.5, 10.3) {寒武纪};
% Habana — acquired by Intel
\fill[figRed, bubble] (6.5, 9.3) circle (6pt);
\node[below, font=\sffamily\tiny] at (6.5, 8.9) {Habana};
% Graphcore — acquired
\fill[figRed, bubble] (2.5, 10.5) circle (5pt);
\node[above, font=\sffamily\tiny] at (2.5, 10.8) {Graphcore};

% ====== Training (y ~ 7.8) ======
% Mosaic ML (Databricks) — acquired
\fill[figRed, bubble] (5, 7.8) circle (8pt);
\node[above, font=\sffamily\tiny] at (5, 8.3) {MosaicML};
% Weights \& Biases — $8.7B
\fill[figBlue, bubble] (6, 7.5) circle (10pt);
\node[below, font=\sffamily\tiny\bfseries] at (6, 7.0) {W\&B};
% Anyscale (Ray) — $1B val
\fill[figBlue, bubble] (4.5, 7.2) circle (7pt);
\node[below left, font=\sffamily\tiny] at (4.2, 6.8) {Anyscale};
% Determined AI — acquired by HPE
\fill[figRed, bubble] (3.2, 7.8) circle (5pt);
\node[above left, font=\sffamily\tiny] at (2.8, 8.1) {Determined AI};
% Lightning AI
\fill[figBlue, bubble] (2.5, 7.5) circle (5pt);
\node[below, font=\sffamily\tiny] at (2.5, 7.1) {Lightning AI};
% Modal
\fill[figBlue, bubble] (3.8, 8.2) circle (6pt);
\node[above, font=\sffamily\tiny] at (3.8, 8.6) {Modal};

% ====== Inference (y ~ 5.8) ======
% Fireworks AI — $2.5B val
\fill[figBlue, bubble] (5, 5.8) circle (8pt);
\node[above, font=\sffamily\tiny] at (5, 6.3) {Fireworks};
% Baseten — $1.1B val
\fill[figBlue, bubble] (4, 5.5) circle (7pt);
\node[below, font=\sffamily\tiny] at (4, 5.0) {Baseten};
% Replicate
\fill[figBlue, bubble] (3.5, 6.2) circle (6pt);
\node[above left, font=\sffamily\tiny] at (3.2, 6.5) {Replicate};
% OctoAI — acquired
\fill[figRed, bubble] (3, 5.5) circle (5pt);
\node[below left, font=\sffamily\tiny] at (2.6, 5.2) {OctoAI};
% vLLM (open source / UC Berkeley)
\fill[figBlue, bubble] (2, 6.0) circle (4pt);
\node[left, font=\sffamily\tiny] at (1.6, 6.0) {vLLM};
% BentoML
\fill[figBlue, bubble] (2.5, 5.8) circle (4pt);
\node[below, font=\sffamily\tiny] at (2.5, 5.4) {BentoML};

% ====== Networking (y ~ 3.8) ======
% Astera Labs — IPO
\fill[figGreen, bubble] (7, 3.8) circle (9pt);
\node[above, font=\sffamily\tiny\bfseries] at (7, 4.3) {Astera Labs};
% Celestial AI
\fill[figBlue, bubble] (2.5, 3.8) circle (6pt);
\node[below, font=\sffamily\tiny] at (2.5, 3.4) {Celestial AI};
% Ayar Labs
\fill[figBlue, bubble] (3.5, 4.1) circle (5pt);
\node[above, font=\sffamily\tiny] at (3.5, 4.5) {Ayar Labs};
% Lightmatter
\fill[figBlue, bubble] (3, 3.5) circle (6pt);
\node[below, font=\sffamily\tiny] at (3, 3.0) {Lightmatter};

% ====== Edge & CDN (y ~ 1.8) ======
% Useful Sensors
\fill[figBlue, bubble] (1.8, 1.8) circle (4pt);
\node[right, font=\sffamily\tiny] at (2.1, 1.8) {Useful Sensors};
% Qualcomm AI (not startup, omit)
% Hailo — $2B
\fill[figBlue, bubble] (3.5, 1.8) circle (6pt);
\node[above, font=\sffamily\tiny] at (3.5, 2.2) {Hailo};
% Kneron
\fill[figBlue, bubble] (2.5, 2.1) circle (4pt);
\node[above, font=\sffamily\tiny] at (2.5, 2.5) {Kneron};
% Mythic — shut down
\fill[figGray, bubble, fill opacity=0.5] (1.5, 2.3) circle (4pt);
\node[left, font=\sffamily\tiny, figGray] at (1.2, 2.3) {Mythic};

\end{tikzpicture}
\caption{基础设施层市场竞争格局全景图。气泡大小近似反映公司估值（截至2026年3月），
横轴为收入规模，纵轴为技术层级。颜色编码：蓝色=独立运营，红色=已被收购，
绿色=已IPO/上市，灰色=已关停或转型。}
\label{fig:infra-market-map}
\end{figure}

% Written 2026-03-21 after cross-section review of §2.1--§2.7.

\subsection{六层价值链与利润分布}

将前七节的公司数据叠加在一起，AI 基础设施层的价值链可沿
六个层级展开：\textbf{GPU 云}（CoreWeave、Lambda、Nebius）、
\textbf{芯片}（NVIDIA + 挑战者）、
\textbf{训练}（Anyscale、MosaicML、Lightning AI）、
\textbf{推理}（Fireworks、Together、Baseten）、
\textbf{网络}（Lightmatter、Ayar Labs、Enfabrica）、
\textbf{边缘}（Hailo、Cloudflare Workers AI、Apple Neural Engine）。
每一层都有数十家创业公司竞争，但利润分布却极度不均。

\begin{keybox}[谁拿走了最多利润？]
NVIDIA 2025 财年数据中心收入超过 \$1,150 亿，
净利润率约 55\%——这意味着仅 NVIDIA 一家，
就从 AI 基础设施层攫取了超过 \$600 亿净利润。
相比之下，CoreWeave 2025 年营收 \$51.3 亿但仍在亏损，
全部 Neocloud 创业公司的营收总和不到 NVIDIA 的十分之一，
且多数处于烧钱扩张阶段。
这形成了一个鲜明的``卖铲子''格局：
\textbf{GPU 云公司是 NVIDIA 的大客户，
芯片挑战者是 NVIDIA 的潜在替代者，
推理公司是 NVIDIA 生态的软件层——
而 NVIDIA 自身则通过硬件销售、软件锁定（CUDA）
和战略收购（Groq、Run:ai、OctoAI、Lepton AI、Enfabrica），
从每一层抽取价值。}
\end{keybox}

网络互联层正在成为被低估的价值节点。
Lightmatter 估值 \$44 亿、Ayar Labs 估值 \$37.5 亿、
Celestial AI 被 Marvell 以 \$32.5 亿收购——
三家光互联公司合计估值超过 \$110 亿，
说明资本已意识到：在万卡集群时代，
\textbf{连接 GPU 的管道与 GPU 本身同等重要}。

\subsection{Neocloud 竞争格局：从百花齐放到寡头整合}

\S\ref{sec:infra-gpu-cloud} 梳理了 11 家 GPU 云创业公司，
但全球 Neocloud 总数已超过 100 家。
这一赛道正经历从``百花齐放''到``寡头整合''的剧烈收缩。

整合逻辑来自三个结构性力量。
\textbf{第一，资本壁垒指数级攀升}。
2023 年建设一个万卡 H100 集群需要约 \$3 亿；
到 2026 年，CoreWeave 单年 capex 已达 \$230 亿，
Nebius 获得 Meta \$270 亿合同——
能匹配这一资本门槛的创业公司屈指可数。
\textbf{第二，NVIDIA 的``Exemplar Cloud''认证}机制
赋予了少数几家公司（CoreWeave、Lambda、Crusoe）
优先获取最新 GPU 的特权，
形成了一种``被选中者才能生存''的马太效应。
\textbf{第三，长期合约锁定}——
CoreWeave 的 \$556 亿积压订单和 Nebius 的 Meta 合约
意味着头部玩家已将未来数年的算力需求预先锁定，
后来者面临的不是技术竞争，而是客户已被瓜分的残酷现实。

\begin{corebox}[两种 Neocloud 生存模式]
\textbf{``超级 Neocloud''模式}（CoreWeave、Nebius、Lambda）：
自建 GW 级数据中心，签下数十亿美元长期合约，
通过纵向整合（CoreWeave 收购 W\&B、Lightning AI 合并 Voltage Park）
向全栈 AI 云演进。5--10 家将存活。

\textbf{``轻资产市场''模式}（RunPod、Vast.ai、Modal、SF Compute）：
不拥有或少量拥有 GPU，以市场撮合、Serverless 抽象
或开发者体验取胜。资本效率极高
（RunPod 融资 \$22M 即达 \$1.2 亿 ARR），
但规模天花板和企业客户信任度是核心瓶颈。
这一层更可能以数十家的形式长期存在，
服务长尾开发者和中小企业。
\end{corebox}

未来两年，Lightning AI/Voltage Park 式的合并将密集发生。
那些既无 NVIDIA 特殊关系、又无差异化定位的中型 Neocloud，
将面临被收购或淘汰的命运。

\subsection{芯片创业的生死博弈}

\S\ref{sec:infra-chips} 覆盖了西方 9 家和中国 10 家 AI 芯片创业公司。
一个残酷的统计是：在西方 9 家中，
Groq 被 NVIDIA \$200 亿收购、Graphcore 被 SoftBank \$4 亿``贱卖''、
Etched 和 MatX 尚未出货——
\textbf{真正以独立公司形态实现规模化产品交付的，
仅有 Cerebras、Tenstorrent 和 Hailo 三家}。

\begin{warnbox}[为什么 90\% 的 AI 芯片创业公司会失败]
AI 芯片创业面临三重死亡螺旋：
\begin{enumerate}[noitemsep]
  \item \textbf{CUDA 生态的引力}。
        十余年积累的数百万开发者、数千个优化库
        构成了从 NVIDIA 迁出的巨大摩擦成本。
        Intel Gaudi 3 定价仅 H100 一半，
        却因软件生态不成熟而出货目标下调 30\%。
  \item \textbf{流片的``死亡之谷''}。
        从芯片设计到量产出货通常需要 3--5 年、
        数亿美元投入，且一次流片失败可能
        直接耗尽公司现金储备。
        Graphcore 从 2016 年成立到 2024 年被收购，
        累计消耗 \$7.67 亿仍未实现可持续商业化。
  \item \textbf{架构押注的赌博性}。
        Etched 和 MatX 押注 Transformer 将主导 5--10 年，
        但 State Space Model 等新架构的崛起
        随时可能让专用 ASIC 变成废硅——
        这正是 Graphcore IPU 在 LLM 时代失去价值的翻版。
\end{enumerate}
\end{warnbox}

中国 AI 芯片的生存逻辑截然不同。
出口管制创造了约 500 亿美元的``真空市场''，
政府 3440 亿元大基金三期和约 \$70B 两会补贴
提供了市场经济中不可能获得的资本支持，
客户``被迫试用''打破了 CUDA 生态锁定。
结果是：寒武纪 2025 年首次盈利（净利 20.59 亿元），
海光信息营收 143.76 亿元，
``四小龙''密集 IPO 且首日涨幅动辄数百个百分点。
\textbf{中国芯片是``政策驱动''的生存逻辑，
西方芯片是``市场驱动''的生死博弈——
前者的风险在于估值泡沫（四家合计市值 1.66 万亿元，PS 超 70 倍），
后者的风险在于独立生存的概率极低。}

\subsection{推理经济学的结构性变革}

\S\ref{sec:infra-inference} 记录了推理赛道最惊人的数据：
\textbf{2022--2026 年间，同等能力的推理成本下降了约 1000 倍}——
从 GPT-3.5 Turbo 的 \$2.00/M tokens
到 DeepSeek-R1 的 \$0.15/M tokens，
小模型更低至 \$0.03/M tokens。

这一降幅的驱动力来自四个叠加层：
PagedAttention 将显存利用率从 30--50\% 提升至 90\%+；
FP8/INT4 量化实现 2--4 倍加速；
Speculative Decoding 进一步压缩延迟；
前缀缓存在 Agent 场景下节省 30--50\% 计算。
这些技术的复合效应，加上 GPU 算力本身的增长，
共同推动了推理成本的指数级下降
（详见 \textit{Emergence} 第九章的技术分析）。

推理优化公司的护城河究竟在哪里？
答案正在从``更便宜的推理''转向``更完整的平台''。
当 vLLM 和 SGLang 作为开源引擎被所有人使用时，
纯粹的性价比竞争必然走向零利润。
Fireworks AI 的``复合 AI 推理引擎''、
Together AI 的``全栈 AI 云''、
Baseten 的企业级推理编排——
它们的共同方向是向上攀升到\textbf{平台层}，
用多模型编排、微调、Agent 工作流等更高层抽象
制造迁移成本。

\begin{keybox}[NVIDIA 系统性收购的防守逻辑]
NVIDIA 在 2024--2025 年连续收购了
Groq（\$200 亿）、Run:ai（\$7 亿）、
OctoAI（$\sim$\$1.65 亿）、Lepton AI（数亿美元）、
Enfabrica（$>$\$9 亿），
并大额投资了 Baseten（\$1.5 亿）和 Together AI。
这并非简单的消灭竞争——而是一套完整的防御体系：
\textbf{控制推理软件栈}（防止 vLLM 在 AMD 上跑得同样好）、
\textbf{获取顶尖人才}（Caffe 之父、TPU 发明者、TVM 核心开发者），
以及\textbf{补全技术拼图}
（Groq 的确定性推理 + 自家 GPU 的训练和 prefill）。
当 AMD ROCm 日益成熟时，
NVIDIA 用软件层的深度绑定制造额外的迁移壁垒。
\end{keybox}

\subsection{2026 年基础设施层展望}

\textbf{谁会成为``AI 版 AWS''？}
CoreWeave 已是最接近的候选者——
\$51.3 亿营收、\$556 亿积压订单、
收购 W\&B 构建全栈平台。
但其客户集中度（OpenAI 占绝对多数）是致命软肋。
Nebius 背靠 Meta \$270 亿合同，在欧洲以 GDPR 合规取胜。
最终，``AI 版 AWS''可能不是某一家 Neocloud，
而是 NVIDIA 自身——通过 DGX Cloud、NIM 和系列收购，
NVIDIA 正在从芯片公司蜕变为 AI 基础设施平台公司。

\textbf{能源将成为终极瓶颈}。
\S\ref{sec:infra-factories} 的数据触目惊心：
五大超算厂商 2026 年 capex 预计达 \$6,350--6,650 亿，
美国数据中心耗电占全国用电 4.4\% 且到 2030 年将翻倍，
科技巨头核电合同总额已超 \$520 亿。
xAI Colossus 目标 2\,GW、Meta Hyperion 规划 5\,GW——
当单一站点用电超过小型国家时，
\textbf{GPU 不再是约束条件，电力才是}。
Crusoe 用废气发电、中东用廉价天然气+光伏——
能获取低成本能源的公司，将拥有结构性成本优势。

\textbf{光互联是否会改变游戏规则？}
以太网在 2025 年首次超越 InfiniBand 占据 AI 后端网络主导份额，
UEC Spec 1.0 标准化了高性能 AI 以太网。
光互联三巨头（Lightmatter、Ayar Labs、Celestial AI/Marvell）
的共封装光学技术正在突破铜缆的带宽和距离极限。
当集群规模从万卡走向百万卡时，
\textbf{光互联不是锦上添花，而是物理必然}——
它将决定下一代 AI 工厂的架构上限。

\begin{corebox}[基础设施层 2026 年五大判断]
\begin{enumerate}[noitemsep]
  \item \textbf{Neocloud 整合加速}：
        100+ 家 GPU 云将在 2026--2027 年缩减至 20--30 家，
        头部 5 家占据 80\%+ 市场份额。
  \item \textbf{芯片格局固化}：
        NVIDIA 维持 80\%+ 训练份额，
        推理芯片赛道留给 Cerebras、Tenstorrent
        和中国国产方案各约 5--10\% 的空间。
  \item \textbf{推理成本继续下降}：
        2026 年底 GPT-4 等效推理将降至 \$0.05/M tokens 以下，
        AI Agent 的大规模部署成为经济可行。
  \item \textbf{能源成为核心竞争力}：
        拥有低成本电力合约（核能/可再生能源）的公司
        将获得 30--50\% 的运营成本优势。
  \item \textbf{光互联标准化}：
        CPO 将从实验走向量产部署，
        成为 10 万+ GPU 集群的标配互联方案。
\end{enumerate}
\end{corebox}

基础设施层是 AI 创业生态的地基。
本章七节覆盖的近百家公司——
从 CoreWeave 的 \$435 亿市值到 JarvisLabs 的印度学生云，
从 NVIDIA 的 \$600 亿净利润到 Banana 的关停倒闭——
共同描绘了一幅残酷而壮观的图景：
\textbf{基础设施层的赢家将定义未来十年 AI 计算的成本曲线，
而成本曲线将决定上层每一个 AI 应用的生死}。
在这个意义上，本章讨论的不仅是芯片和数据中心——
而是整个 AI 时代的地质底层。
